\documentclass{IEEEtran}

\usepackage{graphicx}
\usepackage{algorithm}
\usepackage{amsmath}
\usepackage{amssymb}
\usepackage{cancel}
\usepackage{mathdots}
\usepackage{mathtools}
\usepackage{esint}
\usepackage[version=4]{mhchem}
\usepackage{stackrel}
\usepackage{hyperref}
\hypersetup{
  colorlinks=true,
  linkcolor=blue,
  citecolor=blue,
  urlcolor=blue,
  filecolor=blue
}
% for subfigures/subtables
\usepackage[caption=false, font=footnotesize]{subfig}
\usepackage{algorithmic}
\usepackage{cite}

\newtheorem{assumption}{Assumption}
\newtheorem{lemma}{Lemma}
\newtheorem{theorem}{Theorem}
\newtheorem{definition}{Definition}
\newtheorem{remark}{Remark}

% Loosen line-breaking constraints for the auto-converted manuscript.
\emergencystretch=2em
\hfuzz=20pt
\hbadness=10000
\vbadness=10000

\begin{document}

\title{Constrained Reinforcement Learning with Heterogeneous Knowledge Fusion via Policy Reuse Optimization}
\author{An Liu, Senior Member, IEEE , Zheyuan Zhou and Kexuan Wang \thanks{An Liu and Kexuan Wang are with the College of Information Science and Electronic Engineering, Zhejiang University, Hangzhou 310027, China (email: \{anliu, kexuanWang\}@zju.edu.cn). (Corresponding Author: Zheyuan Zhou)}}
\maketitle
\begin{abstract}
Constrained Reinforcement Learning (CRL) faces significant challenges in practical applications due to its slow convergence and high interaction costs. While both transfer learning (TL) and domain-specific knowledge (DK) offer potential solutions, existing methods often rely on heuristics and lack a unified framework to integrate these heterogeneous knowledge sources with theoretical guarantees. This paper proposes a novel Constrained Policy Reuse Optimization (CPRO) framework that systematically fuses policies from TL (referred to as source policies) and DK-based policies to accelerate and robustify CRL. Within this framework, the fusion of heterogeneous knowledge is formulated as a Constrained Markov Decision Process (CMDP), which jointly optimizes the policy reuse probabilities and a new learnable policy to maximize cumulative reward while satisfying safety constraints. To solve this problem, we introduce the Fused-CPRO algorithm, which leverages an efficient single-loop deep actor-critic (DAC) structure. In the actor module, the constrained stochastic successive convex approximation (CSSCA) method handles non-convex stochastic constraints; in the critic module, Temporal-Difference (TD) learning updates the critic DNNs with only one or a few finite steps per iteration, simplifying the architecture to a single-loop scheme. The algorithm further improves sampling efficiency by reusing both offline data from source policies and outdated online data, thereby minimizing interaction overhead. We provide a rigorous theoretical analysis, proving that Fused-CPRO converges almost surely to a Karush-Kuhn-Tucker (KKT) point and delivers performance superior to all prior policies. Simulations demonstrate that the proposed method achieves superior performance with faster convergence and significantly lower interaction costs compared to state-of-the-art baselines.
\end{abstract}
\begin{IEEEkeywords}
Constrained reinforcement learning, policy reuse, Transfer Learning, domain knowledge.
\end{IEEEkeywords}
\section{Introduction}
Reinforcement learning (RL) has emerged as a powerful paradigm for solving sequential decision-making problems through interactions with dynamic environments, with notable applications ranging from robot navigation \cite{robotlocomotion} to radio resource management in wireless communications \cite{RRM}. However, the deployment of conventional reinforcement learning (RL), particularly deep RL (DRL), in practical systems like wireless networks remains challenging. Such environments are often characterized by non-stationarity and stringent safety requirements, which expose two fundamental limitations of existing methods. First, they suffer from sample inefficiency, requiring an enormous number of costly interactions with the environment to converge. Second, they exhibit limited robustness: the black-box nature of deep neural networks (DNNs) can lead to unreliable and unpredictable behavior when the agent encounters out-of-distribution states, posing significant safety risks.

The challenge of sample inefficiency is commonly addressed through transfer learning (TL). The core idea of TL in RL is to repurpose knowledge--- such as learned policies, historical data, or favorable initializations--- from related tasks to accelerate learning and improve initial performance in a target task. In the broader RL literature, closely related mechanisms include policy reuse and approximate policy iteration style updates \cite{boundedJ1}, off-policy or offline data reuse to reduce fresh interactions with the environment \cite{SCAOPO,offlineCRLrecommendedbyreviewer2}, warm-starting deep models with better parameter initialization \cite{initialization}, and prior-guided reward or objective design as discussed in standard RL references \cite{RLintroduction,DRL}.

In parallel to TL, another line of research focuses on infusing explicit prior structure--- especially safety constraints and cautious exploration mechanisms--- directly into the RL process \cite{CPO,CLQset,reviewer1_1_3}. The primary goal of incorporating such domain knowledge is to enhance the robustness, safety, and interpretability of Deep RL agents, which are often criticized as black-box systems. A common technique is to bias exploration toward safe regions or apply responsive safety corrections, allowing the agent to avoid catastrophic failures or clearly suboptimal decisions during learning \cite{SCAOPO17,SCAOPO18,PPOLagTRPOLag}. While TL aims to improve data efficiency by leveraging past experience, DK seeks to improve reliability by injecting prior understanding of the domain.

Despite these advances, two critical limitations persist. First, most existing TL and DK methods are heuristic in design, lacking rigorous theoretical guarantees for convergence or performance. This shortcoming is most apparent in policy reuse, where selection probabilities are typically set in an ad-hoc manner rather than being systematically optimized. Second, there is no unified framework to systematically integrate TL and DK, despite being complementary sources of heterogeneous knowledge. This gap is especially pronounced in Constrained RL (CRL) settings, where agents must satisfy complex (often non-convex and stochastic) safety or cost constraints during the learning process itself. To bridge this gap, we propose a Constrained Policy Reuse Optimization (CPRO) framework that formally unifies TL and DK through a constrained Markov decision process (CMDP). Within this formulation, we jointly optimize a new learnable policy and reuse probabilities for both TL-sourced and DK-derived policies, enabling dynamic, performance-driven knowledge fusion. Our main contributions are:

\begin{itemize}
\item A unified fusion framework for CRL:

We introduce the CPRO framework, which integrates heterogeneous knowledge sources, i.e., TL policies and DK policies, into a principled CMDP. The framework automatically adjusts reuse probabilities to accelerate learning and enhance robustness, falling back to DK policies when DNN-based policies underperform.
\item An efficient single-loop deep actor-critic (DAC) algorithm:

We develop Fused-CPRO, a deep actor-critic algorithm that combines constrained stochastic successive convex approximation (CSSCA) in the actor with lightweight TD updates in the critic. This single-loop design--- enabled by finite-step or even singl-step critic updates--- improves efficiency while supporting mixed offline-online data reuse from both source policies and past interactions.
\item Theoretical guarantees:

We prove that Fused-CPRO converges almost surely to a KKT point of the CPRO problem and performs no worse than the best prior policy, with the potential to surpass all existing policies under mild representational conditions.
\end{itemize}
The remainder of the paper is organized as follows: Section II formulates the constrained policy reuse optimization problem. Sections III and IV present the Fused-CPRO algorithm and its theoretical analysis, respectively. Section V provides simulation results, and Section VI concludes the paper.

\section{Problem Formulation for Constrained Policy Reuse Optimization}
\label{sec:Preliminaries}
This section establishes the formal framework for our work. We first review the fundamentals of CMDP. We then introduce the novel constrained policy reuse mechanism designed to integrate heterogeneous knowledge from TL and DK. Finally, we formulate the joint optimization of policy reuse probabilities and a new learnable policy as a CMDP problem.

\subsection{Preliminaries on CMDP}
\label{subsec:Problem-Formulation}
A CMDP is defined by the tuple $(\mathcal{S},\mathcal{A},P,\mathcal{C})$:

\begin{itemize}
\item $\mathcal{S}\subseteq\mathbb{R}^{n_{s}}$: The state space.
\item $\mathcal{A}\subseteq\mathbb{R}^{n_{a}}$: The action space.
\item $P:\mathcal{S}\times\mathcal{A}\times\mathcal{S}\rightarrow[0,1]$: The state transition probability function, where $P(\boldsymbol{s}'\mid\boldsymbol{s},\boldsymbol{a})$ denotes the probability of transitioning to state $\boldsymbol{s}'$ upon taking action $\boldsymbol{a}$ in state $\boldsymbol{s}$.
\item $\mathcal{C}=\{C_{i}\}_{i=0}^{I}$: A set of per-stage cost functions. Typically,

$C_{0}$ represents the primary cost to be minimized (often the negative reward), while $\{C_{i}\}_{i=1}^{I}$ represent auxiliary costs to be constrained.
\end{itemize}

A policy $\pi:\mathcal{S}\rightarrow\mathcal{P}(\mathcal{A})$ maps states to probability distributions over actions, with $\pi(\boldsymbol{a}\mid\boldsymbol{s})$ being the probability of selecting action $\boldsymbol{a}$ in state $\boldsymbol{s}$. The policy $\pi$ and the transition dynamics $P$ together induce a probability distribution over trajectories $\tau=\{\boldsymbol{s}_{0},\boldsymbol{a}_{0},\boldsymbol{s}_{1},\boldsymbol{a}_{1},\ldots\}$, denoted by $p_{\pi}$, such that $\boldsymbol{s}_{t+1}\sim P(\cdot\mid\boldsymbol{s}_{t},\boldsymbol{a}_{t})$ and $\boldsymbol{a}_{t}\sim\pi(\cdot\mid\boldsymbol{s}_{t})$.

The standard infinite-horizon average-cost CMDP problem is formulated as:
\begin{align}
\underset{\pi}{\mathrm{min}}J_{0}\left(\boldsymbol{\theta}\right)\overset{\triangle}{=} & \underset{T\rightarrow\infty}{\mathrm{lim}}\frac{1}{T}\mathbb{E}_{p_{\pi}}\left[\stackrel[t=0]{T-1}{\sum}C_{0}\left(\boldsymbol{s}_{t},\boldsymbol{a}_{t}\right)\right]\label{eq:problem 1}\\
\mathrm{s.t.}J_{i}\left(\boldsymbol{\theta}\right)\overset{\triangle}{=} & \underset{T\rightarrow\infty}{\mathrm{lim}}\frac{1}{T}\mathbb{E}_{p_{\pi}}\left[\stackrel[t=0]{T-1}{\sum}C_{i}\left(\boldsymbol{s}_{t},\boldsymbol{a}_{t}\right)\right]-c_{i}\leq0,\nonumber \\
& i=1,\ldots,I,\nonumber
\end{align}
where $c_{i}\in\mathbb{R}$ are predefined constraint thresholds.

\subsection{Policy Reuse for Heterogeneous Knowledge Integration}
In policy reuse, there are $N\geq1$ old policies $\pi_{1},...,\pi_{N}$ and a new policy $\pi_{0}$. At each time $t$, the $n$-th policy is used with probability $\rho_{n}$. In general, there are two subsets of old policies. The first subset $\pi_{n},n=1,...,N_{1}$ contains $N_{1}$ old policies obtained from TL, which are called TL policies. The second subset $\pi_{n},n=N_{1}+1,...,N$ contains $N_{2}=N-N_{1}$ old policies constructed using domain knowledge, which are called DK policies. The DK policies are often an iterative algorithm designed using the domain knowledge. On the other hand, the TL polices and new policy are parameterized with DNNs, and we use $\boldsymbol{\psi}_{n}\in\boldsymbol{\Psi}$ to denote the parameters of the DNN for the $n$-th policy with $n=0,1,...,N_{1}$. The old policies are fixed and the new policy together with the reuse probabilities $\boldsymbol{\rho}=\left[\rho_{0},\rho_{1},...,\rho_{n}\right]^{T}$ will be updated using SCADAC-PRO.

Our framework employs a policy reuse mechanism that leverages $N\geq1$ old policies $\{\pi_{n}\}_{n=1}^{N}$ alongside a new, learnable policy $\pi_{0}$. At each timestep $t$, the agent selects policy $\pi_{n}$ with probability $\rho_{n}$ to generate the action. The set of old policies comprises two distinct categories of heterogeneous knowledge:

\begin{enumerate}
\item Transfer Learning (TL) Policies:

The first $N_{1}$ policies, $\{\pi_{n}\}_{n=1}^{N_{1}}$, are TL policies acquired from related source tasks. These are typically parameterized by DNNs with parameters $\boldsymbol{\psi}_{n}\in\boldsymbol{\Psi}$.
\item Domain Knowledge (DK) Policies:

The remaining $N_{2}=N-N_{1}$ policies, $\{\pi_{n}\}_{n=N_{1}+1}^{N}$, are DK policies. These are often deterministic algorithms or rule-based controllers constructed from prior expert knowledge about the problem domain.
\end{enumerate}

The new policy $\pi_{0}$ is also DNN-parameterized with parameters $\boldsymbol{\psi}_{0}$. During optimization, all old policies remain fixed. The learning variables are the reuse probability vector $\boldsymbol{\rho}=[\rho_{0},\rho_{1},\dots,\rho_{N}]^{T}$ and the new policy's parameters $\boldsymbol{\psi}_{0}$, which are updated jointly by our proposed algorithm. To enable gradient-based optimization and ensure a unified representation, we model the TL policies and the new policy as Gaussian policies \cite{Gaussianpolicy}. For policy $n\in\{0,1,\ldots,N_{1}\}$, the mean $\boldsymbol{\mu}_{n}$ and the diagonal elements of the covariance matrix $\boldsymbol{\Sigma}_{n}$ are output by DNNs $f_{m_{\boldsymbol{\mu}}}(\boldsymbol{\psi}_{n};\boldsymbol{s})$ and $f_{m_{\boldsymbol{\sigma}}}(\boldsymbol{\psi}_{n};\boldsymbol{s})$, respectively. We restrict $\boldsymbol{\Sigma}_{n}$ to be a diagonal matrix (i.e., its off-diagonal elements are zero), a common practice that reduces the number of parameters and stabilizes learning while retaining sufficient expressiveness for exploration. Thus, the policy is defined as:
\[
\pi_{n}(\boldsymbol{a}\mid\boldsymbol{s})\propto|\boldsymbol{\Sigma}_{n}|^{-\frac{1}{2}}\exp\left(-\frac{1}{2}(\boldsymbol{\mu}_{n}-\boldsymbol{a})^{\intercal}\boldsymbol{\Sigma}_{n}^{-1}(\boldsymbol{\mu}_{n}-\boldsymbol{a})\right).
\]

DK policies, however, are often deterministic (e.g., rule-based controllers or iterative algorithms). To incorporate them into our stochastic policy gradient framework, we apply a Gaussian smoothing technique. Specifically, each DK policy is treated as a Gaussian $\mathcal{N}(\boldsymbol{\mu}_{n},\epsilon\mathbf{I})$, where $\boldsymbol{\mu}_{n}$ is its deterministic output and $\epsilon$ is a small, fixed variance. This approximation serves a dual purpose: it provides a differentiable surrogate for gradient computation through the probability density, while the small variance ensures the sampled actions remain close to the original deterministic output, preserving the intended behavior of the DK policy.

\subsection{Formulation of the Constrained Policy Reuse Optimization Problem}
The described policy reuse scheme defines a mixed policy $\pi_{\boldsymbol{\theta}}$, parameterized by $\boldsymbol{\theta}=[\boldsymbol{\rho};\boldsymbol{\psi}_{0}]\in\boldsymbol{\Theta}$. At each step, an action is sampled by first selecting a policy index $n$ according to $\boldsymbol{\rho}$ and then sampling an action from $\pi_{n}$. The joint optimization of the reuse probabilities $\boldsymbol{\rho}$ and the new policy parameters $\boldsymbol{\psi}_{0}$ is thus formulated as the following CMDP:
\begin{align}
\underset{\boldsymbol{\theta}\in\boldsymbol{\Theta}}{\mathrm{min}}J_{0}\left(\boldsymbol{\theta}\right)\overset{\triangle}{=} & \underset{T\rightarrow\infty}{\mathrm{lim}}\frac{1}{T}\mathbb{E}_{p_{\pi_{\boldsymbol{\theta}}}}\left[\stackrel[t=0]{T-1}{\sum}C_{0}\left(\boldsymbol{s}_{t},\boldsymbol{a}_{t}\right)\right]\label{eq:problem 1-1}\\
\mathrm{s.t.}J_{i}\left(\boldsymbol{\theta}\right)\overset{\triangle}{=} & \underset{T\rightarrow\infty}{\mathrm{lim}}\frac{1}{T}\mathbb{E}_{p_{\pi_{\boldsymbol{\theta}}}}\left[\stackrel[t=0]{T-1}{\sum}C_{i}\left(\boldsymbol{s}_{t},\boldsymbol{a}_{t}\right)\right]-c_{i}\leq0,\nonumber \\
& i=1,\ldots,I,\nonumber
\end{align}
where $p_{\pi_{\boldsymbol{\theta}}}$ denotes the trajectory distribution under the mixed policy $\pi_{\boldsymbol{\theta}}$.

In the above CMDP formulation, we focus on the case of infinite-horizon average reward/cost. The extension to discounted case is strightfoward \cite{reviewer1_4_3}.

\section{Fused-CPRO Algorithm}
\label{sec:Algorithmic-Framework}
\begin{figure}[t]
\noindent
\centering

\includegraphics[width=8cm,height=6cm]{algorithm.png}

\caption{The algorithmic framework of the proposed SLDAC}\label{fig:diagram of the algorithm}
\end{figure}

This section details the proposed Fused Policy Reuse Optimization (Fused-CPRO) algorithm. As illustrated in Fig. \ref{fig:diagram of the algorithm}, the algorithm operates in a single-loop fashion, iterating between a critic module for state-action value function (i.e., Q-function) estimation and an actor module for policy optimization. At each iteration $t$, a new online data sample $\varepsilon_{t}$ is collected into a replay buffer. The updates for both the critic and the actor then leverage a mixed offline-online data reuse scheme to enhance sample efficiency. We begin by describing the data collection strategy, followed by the derivation of the policy reuse gradient, and finally the detailed designs of the critic and actor modules. The complete algorithm is summarized in Algorithm \ref{alg:the-Single-Loop-Deep}.

\subsection{Mixed Offline-Online Data Collection}
A key feature of Fused-CPRO is its efficient use of historical data to overcome sample inefficiency. This is achieved through a mixed offline-online data reuse scheme, which strategically leverages two distinct data sources to accelerate learning and reduce fresh interactions with the environment. Specifically, defining the adjusted cost functions as $C_{0}^{\text{'}}(\boldsymbol{s},\boldsymbol{a})=C_{0}(\boldsymbol{s},\boldsymbol{a})$ and $C_{i}^{\text{'}}(\boldsymbol{s},\boldsymbol{a})=C_{i}(\boldsymbol{s},\boldsymbol{a})-c_{i}$ for $i=1,\ldots,I$, then the two data sources are:

\begin{itemize}
\item Offline Datasets ($\mathcal{D}_{n}$):

The agent has access to $N$ pre-collected datasets. Each dataset $\mathcal{D}_{n}$ contains tuples $\left\{ \boldsymbol{s},\boldsymbol{a},\{C'_{i}(\boldsymbol{s},\boldsymbol{a})\}_{i=0}^{I},\boldsymbol{s}'\right\} $ generated by executing the $n$-th old policy $\pi_{n}$ in an environment similar to the current target task.
\item Online Dataset ($\mathcal{D}_{0}^{t}$):

Concurrently, the agent interacts with the current environment using the latest mixed policy $\pi_{\boldsymbol{\theta}_{t}}$. At iteration $t$, it obtains a new sample $\varepsilon_{t}=\{\boldsymbol{s}_{t},\boldsymbol{a}_{t},\{C_{i}(\boldsymbol{s}_{t},\boldsymbol{a}_{t})\}_{i=0}^{I},\boldsymbol{s}_{t+1}\}$ and maintains a replay buffer $\mathcal{D}_{0}^{t}$ containing the latest $T_{t}$ samples: $\mathcal{D}_{0}^{t}=\{\varepsilon_{t-T_{t}+1},...,\varepsilon_{t}\}$.
\end{itemize}
\begin{remark}
In practice, a mini-batch of $B$ samples can be collected to perform $q$ critic updates per iteration ($B\geq q\geq1$), allowing a flexible trade-off between gradient estimation accuracy and interaction cost. For clarity in algorithm presentation and convergence analysis, we focus on the case of one sample per iteration ($B=q=1$). The theoretical guarantees in Section \ref{sec:CONVERGENCE-ANALYSIS} hold for more general case with $B\geq q\geq1$.

\end{remark}

\subsection{Policy Reuse Gradient with Mixed Data}
Let $\pi_{\boldsymbol{\psi}_{0}}=\pi_{0}$ denote the new, learnable policy. The mixed policy is a weighted sum: $\pi_{\boldsymbol{\theta}}(\boldsymbol{a}|\boldsymbol{s})=\sum_{n=0}^{N}\rho_{n}\pi_{n}(\boldsymbol{a}|\boldsymbol{s})$, with $\boldsymbol{\theta}=[\boldsymbol{\rho};\boldsymbol{\psi}_{0}]$. The goal is to find the gradient of the objective and constraint functions $J_{i}(\boldsymbol{\theta})$ with respect to $\boldsymbol{\theta}$.

Following the policy gradient theorem, the gradient can be expressed as:

\[
\nabla_{\boldsymbol{\theta}}J_{i}(\boldsymbol{\theta})=\mathbb{E}_{\sigma_{\pi_{\boldsymbol{\theta}}}}\left[Q_{i}^{\pi_{\boldsymbol{\theta}}}(\boldsymbol{s},\boldsymbol{a})\nabla_{\boldsymbol{\theta}}\log\pi_{\boldsymbol{\theta}}(\boldsymbol{a}|\boldsymbol{s})\right],
\]
where $\sigma_{\pi_{\boldsymbol{\theta}}}$ is the stationary state-action distribution and $Q_{i}^{\pi_{\boldsymbol{\theta}}}$ is the state-action value function. Applying this to our mixed policy yields the gradients for the reuse probabilities and the new policy parameters:

\[
\begin{aligned}\nabla_{\rho_{n}}J_{i}(\boldsymbol{\theta}) & =\mathbb{E}_{\sigma_{\pi_{\boldsymbol{\theta}}}}\left[\frac{Q_{i}^{\pi_{\boldsymbol{\theta}}}(\boldsymbol{s},\boldsymbol{a})\pi_{n}(\boldsymbol{a}|\boldsymbol{s})}{\sum_{k=0}^{N}\rho_{k}\pi_{k}(\boldsymbol{a}|\boldsymbol{s})}\right],\forall n,\\
\nabla_{\boldsymbol{\psi}_{0}}J_{i}(\boldsymbol{\theta}) & =\mathbb{E}_{\sigma_{\pi_{\boldsymbol{\theta}}}}\left[\frac{Q_{i}^{\pi_{\boldsymbol{\theta}}}(\boldsymbol{s},\boldsymbol{a})\rho_{0}\nabla_{\boldsymbol{\psi}_{0}}\pi_{\boldsymbol{\psi}_{0}}(\boldsymbol{a}|\boldsymbol{s})}{\sum_{k=0}^{N}\rho_{k}\pi_{k}(\boldsymbol{a}|\boldsymbol{s})}\right].
\end{aligned}
\]

The key challenge is that the true Q-functions and the expectation under $\sigma_{\pi_{\boldsymbol{\theta}}}$ are unknown. We address this with two approximations:

\begin{enumerate}
\item Function Approximation:

Dual critic DNNs $f(\boldsymbol{\omega}^{i};\boldsymbol{s},\boldsymbol{a})$ and $f(\bar{\boldsymbol{\omega}}^{i};\boldsymbol{s},\boldsymbol{a})$ are used to approximate each $Q_{i}^{\pi_{\boldsymbol{\theta}}}(\boldsymbol{s},\boldsymbol{a})$. Their update is detailed in Section 3.3.
\item Sample Average Approximation:

The expectation is estimated using samples from mixed datasets.
\end{enumerate}

This leads to the mixed-data gradient estimator at iteration $t$:

\begin{equation}
\tilde{\boldsymbol{g}}_{i}^{t}=\left[\tilde{g}_{i,\rho_{1}}^{t};...;\tilde{g}_{i,\rho_{N}}^{t},\tilde{\boldsymbol{g}}_{i,\boldsymbol{\psi}_{0}}^{t}\right]\label{eq:gtlt}
\end{equation}
with the gradient component $\tilde{g}_{i,\rho_{n}}^{t}$ for each reuse probability $\rho_{n}$ and $\tilde{\boldsymbol{g}}_{i,\boldsymbol{\psi}_{0}}^{t}$ for the new policy parameters $\boldsymbol{\psi}_{0}$ respectively given by
\begin{align}
\tilde{g}_{i,\rho_{n}}^{t}= & \xi_{t}\hat{\mathbb{E}}_{\mathcal{D}_{\text{off}}^{t}}\left[\frac{f\left(\bar{\boldsymbol{\omega}}_{t}^{i};\boldsymbol{s},\boldsymbol{a}\right)\pi_{n}\left(\boldsymbol{a}\mid\boldsymbol{s}\right)}{\sum_{k=0}^{N}\rho_{k}\pi_{k}\left(\boldsymbol{a}\mid\boldsymbol{s}\right)}\right]\label{eq:g-tilde-1}\\
& +\left(1-\xi_{t}\right)\hat{\mathbb{E}}_{\mathcal{D}_{0}^{t}}\left[\frac{f\left(\bar{\boldsymbol{\omega}}_{t}^{i};\boldsymbol{s},\boldsymbol{a}\right)\pi_{n}\left(\boldsymbol{a}\mid\boldsymbol{s}\right)}{\sum_{k=0}^{N}\rho_{k}\pi_{k}\left(\boldsymbol{a}\mid\boldsymbol{s}\right)}\right],\forall i,\nonumber
\end{align}

\begin{align}
\tilde{\boldsymbol{g}}_{i,\boldsymbol{\psi}_{0}}^{t}= & \xi_{t}\hat{\mathbb{E}}_{\mathcal{D}_{\text{off}}^{t}}\left[\frac{f\left(\bar{\boldsymbol{\omega}}_{t}^{i};\boldsymbol{s},\boldsymbol{a}\right)\rho_{0}\nabla_{\boldsymbol{\psi}_{0}}\pi_{\boldsymbol{\psi}_{0}}\left(\boldsymbol{a}\mid\boldsymbol{s}\right)}{\sum_{k=0}^{N}\rho_{k}\pi_{k}\left(\boldsymbol{a}\mid\boldsymbol{s}\right)}\right]\label{eq:g-tilde}\\
& +\left(1-\xi_{t}\right)\hat{\mathbb{E}}_{\mathcal{D}_{0}^{t}}\left[\frac{f\left(\bar{\boldsymbol{\omega}}_{t}^{i};\boldsymbol{s},\boldsymbol{a}\right)\rho_{0}\nabla_{\boldsymbol{\psi}_{0}}\pi_{\boldsymbol{\psi}_{0}}\left(\boldsymbol{a}\mid\boldsymbol{s}\right)}{\sum_{k=0}^{N}\rho_{k}\pi_{k}\left(\boldsymbol{a}\mid\boldsymbol{s}\right)}\right],\forall i,\nonumber
\end{align}
Here, $\hat{\mathbb{E}}_{\mathcal{D}}$ denotes the sample mean over dataset $\mathcal{D}$. The mixed offline dataset $\mathcal{D}_{\text{off}}^{t}$ is constructed by sampling $T_{\text{off}}^{t}$ transitions, where each sample is drawn from $\mathcal{D}_{n}$ with probability $\rho_{n}/\sum_{k=1}^{N}\rho_{k}$. The sequence $\{\xi_{t}\}$ is a decreasing weight that balances the contribution of offline and online data. Correspondingly, the function value $J_{i}(\boldsymbol{\theta}_{t})$ is estimated using the same mixture as
\begin{align}
\tilde{J}_{i}^{t} & =\xi_{t}\hat{\mathbb{E}}_{\mathcal{D}_{\text{off}}^{t}}\left[C_{i}^{\text{'}}(\boldsymbol{s},\boldsymbol{a})\right]+\left(1-\xi_{t}\right)\hat{\mathbb{E}}_{\mathcal{D}_{0}^{t}}\left[C_{i}^{\text{'}}(\boldsymbol{s},\boldsymbol{a})\right],\forall i.\label{eq:J-new}
\end{align}

\begin{remark}
This gradient estimator is a key to improve the sample efficiency. It reuses offline data from old policies to warm-start learning and outdated online data (from recent iterations) to decrease gradient variance, significantly accelerating initial convergence without extra interactions.

\end{remark}

\subsection{Critic Module: Learning Surrogate Q-Functions with Stabilized TD Updates}
The critic module is responsible for learning accurate estimates of the state-action value functions (Q-functions), which are essential for computing the policy gradients in (\ref{eq:gtlt}). Since the true average costs $J_{i}(\boldsymbol{\theta}_{t})$ are unknown during learning, we first define a set of surrogate Q-functions $\{\hat{Q}_{i}^{\pi_{\boldsymbol{\theta}_{t}}}\}_{i=0}^{I}$ that use the running estimate $\hat{J}_{i}^{t}$ to replace $J_{i}(\boldsymbol{\theta}_{t})$:
\begin{align}
\hat{Q}_{i}^{\pi_{\boldsymbol{\theta}_{t}}}\left(\boldsymbol{s},\boldsymbol{a}\right)= & \mathbb{E}_{p_{t}}\Bigl[\sum_{l=0}^{\infty}\Bigl(C_{i}^{\text{'}}(\boldsymbol{s}_{l},\boldsymbol{a}_{l})-\hat{J}_{i}^{t}\Bigr)\nonumber \\
& \bigl|\boldsymbol{s}_{0}=\boldsymbol{s},\boldsymbol{a}_{0}=\boldsymbol{a}\Bigr],\forall i,\boldsymbol{s},\boldsymbol{a},\label{eq:definition of Qhat}
\end{align}
where the running estimate $\hat{J}_{i}^{t}$ satisfies $\lim_{t\rightarrow\infty}\bigl|\hat{J}_{i}^{t}-J_{i}\bigl(\boldsymbol{\theta}_{t}\bigr)\bigr|=0$ (see (\ref{eq:J-average}) for the detailed expression), and $p_{t}\triangleq p_{\pi_{\boldsymbol{\theta}_{t}}}$ denotes the trajectory distribution under the current policy. The surrogate $\hat{Q}_{i}^{\pi_{\boldsymbol{\theta}_{t}}}$ converges to the true Q-function $Q_{i}^{\pi_{\boldsymbol{\theta}_{t}}}$ as $t\rightarrow\infty$. To approximate these surrogate Q-functions, we employ dual DNNs, a design that enhances learning stability and is critical for the convergence of actor-critic methods.

\subsubsection{Primary Critic Networks: Minimizing the Mean-Squared Bellman Error}
The first set of DNNs, $\{f(\boldsymbol{\omega}^{i};\boldsymbol{s},\boldsymbol{a})\}_{i=0}^{I}$, are trained to approximate $\hat{Q}_{i}^{\pi_{\boldsymbol{\theta}_{t}}}$. Their parameters $\boldsymbol{\omega}^{i}$ are updated by minimizing the Mean-Squared Bellman Error (MSBE) as
\begin{equation}
\underset{\boldsymbol{\omega}_{t}^{i}}{\mathrm{min}}\ \mathbb{E}_{\sigma_{t}}\Bigl[\Bigl(f\left(\boldsymbol{\omega}_{t}^{i};\boldsymbol{s}_{t},\boldsymbol{a}_{t}\right)-\mathcal{T}_{t}\ f\left(\boldsymbol{\omega}_{t}^{i};\boldsymbol{s}_{t},\boldsymbol{a}_{t}\right)\Bigr)^{2}\Bigr],\forall i,\label{eq:MSBE}
\end{equation}
where $\sigma_{t}\triangleq\sigma_{\pi_{\boldsymbol{\theta}_{t}}}$ is the stationary state-action distribution, and $\mathcal{T}_{t}$ is the Bellman operator for the current policy defined as
\begin{align}
\mathcal{T}_{t}\ f\left(\boldsymbol{\omega}_{t}^{i};\boldsymbol{s}_{t},\boldsymbol{a}_{t}\right)= & \mathbb{E}_{p_{t}}\Bigl[f\left(\boldsymbol{\omega}_{t}^{i};\boldsymbol{s}_{t+1},\boldsymbol{a}_{t+1}'\right)\Bigr]\\
& +C_{i}^{\text{'}}(\boldsymbol{s}_{t},\boldsymbol{a}_{t})-\hat{J}_{i}^{t},\forall i,\nonumber
\end{align}
Here, $\boldsymbol{a}_{t+1}'$ is the action selected by $\pi_{\boldsymbol{\theta}_{t}}$ at the next state $\boldsymbol{s}_{t+1}$. To ensure stable learning, we constrain the DNN parameters within a bounded region. Specifically, we define a ball
\[
\mathbb{B}(\boldsymbol{\omega}_{0}^{i},R_{\boldsymbol{\omega}})=\{\boldsymbol{\omega}^{i}:\|\boldsymbol{\omega}^{i}-\boldsymbol{\omega}_{0}^{i}\|\leq R_{\boldsymbol{\omega}}\}
\]
around the random initialization $\boldsymbol{\omega}_{0}^{i}$. The parameters are then updated via projected TD learning, i.e.,
\begin{equation}
\boldsymbol{\omega}_{t}^{i}=\Pi_{\Omega_{i}}\bigl(\boldsymbol{\omega}_{t-1}^{i}-\eta_{t}\boldsymbol{\Delta}_{i}^{\boldsymbol{\omega}_{t-1}^{i}}\bigr),\forall i,\label{eq:TD-learning}
\end{equation}
where $\{\eta_{t}\}$ is a decreasing step-size sequence satisfying Assumption 2 in section \ref{assumption:step_size}, $\Pi_{\Omega_{i}}$ projects the parameter onto the constraint set $\Omega_{i}\triangleq\mathbb{B}(\boldsymbol{\omega}_{0}^{i},R_{\boldsymbol{\omega}})$, and the TD-error gradient $\boldsymbol{\Delta}_{i}^{\boldsymbol{\omega}}$ is given by
\begin{align}
\boldsymbol{\Delta}_{i}^{\boldsymbol{\omega}_{t-1}^{i}}= & \Bigl(f\left(\boldsymbol{\omega}_{t-1}^{i};\boldsymbol{s}_{t},\boldsymbol{a}_{t}\right)-\bigl(C_{i}^{\text{'}}(\boldsymbol{s}_{t},\boldsymbol{a}_{t})-\hat{J}_{i}^{t-1}\label{eq:gradient term}\\
& +f\left(\boldsymbol{\omega}_{t-1}^{i};\boldsymbol{s}_{t+1},\boldsymbol{a}_{t+1}'\right)\bigl)\Bigr)\nabla_{\boldsymbol{\omega}}f\left(\boldsymbol{\omega}_{t-1}^{i};\boldsymbol{s}_{t},\boldsymbol{a}_{t}\right).\nonumber
\end{align}

\subsubsection{Target Critic Networks: Stabilizing Policy Gradient Estimates}
Directly using the primary critics $f(\boldsymbol{\omega}_{t}^{i})$ in the actor's gradient (\ref{eq:gtlt}) can lead to high variance and instability due to frequent updates. To mitigate this, we employ a second set of target networks $\{f(\bar{\boldsymbol{\omega}}^{i};\boldsymbol{s},\boldsymbol{a})\}_{i=0}^{I}$. Their parameters $\bar{\boldsymbol{\omega}}^{i}$ are updated via a slow-moving average of the primary network parameters, which smooths the target values:
\begin{equation}
\bar{\boldsymbol{\omega}}_{t}^{i}=\left(1-\gamma_{t}\right)\bar{\boldsymbol{\omega}}_{t-1}^{i}+\gamma_{t}\boldsymbol{\omega}_{t}^{i},\forall i,\label{eq:w^_}
\end{equation}
initialized with $\bar{\boldsymbol{\omega}}_{0}^{i}=\boldsymbol{\omega}_{1}^{i}$. Here, $\{\gamma_{t}\}$ is another decreasing sequence. The outputs of these stable target networks, $f(\bar{\boldsymbol{\omega}}_{t}^{i};\boldsymbol{s},\boldsymbol{a})$, are used to compute the policy gradient estimates in (\ref{eq:gtlt}), leading to more reliable and consistent updates for the actor module.

Key Design Rationale: This dual-network architecture, with primary networks trained via TD learning and target networks updated via slow-moving average, is a proven technique to decouple the value estimation and policy improvement processes, significantly enhancing the overall stability and convergence properties of deep actor-critic algorithms.

\subsection{Actor Module: Constrained Optimization via the CSSCA Framework}
The actor module is responsible for updating the policy parameters $\boldsymbol{\theta}=[\boldsymbol{\rho};\boldsymbol{\psi}_{0}]$ by solving the constrained optimization problem in (\ref{eq:problem 1}). Directly optimizing the original non-convex and stochastic objective $J_{0}(\boldsymbol{\theta})$ subject to constraints $J_{i}(\boldsymbol{\theta})\leq0$ is intractable due to the lack of closed-form expressions and the presence of noise in gradient estimates $\tilde{\boldsymbol{g}}_{i}^{t}$. To overcome these challenges, we adopt the CSSCA framework. The core idea is to iteratively replace the original complex problem with a sequence of simpler, deterministic, and convex surrogate problems that are constructed using current stochastic estimates.

The CSSCA procedure in each iteration $t$ consists of three key steps: (1) constructing convex surrogate functions, (2) solving the surrogate optimization problem, and (3) updating the policy parameters.

Step 1: Construction of Convex Surrogate Functions At iteration $t$, we construct a convex quadratic surrogate function $\bar{J}_{i}^{t}(\boldsymbol{\theta})$ that locally approximates each original function $J_{i}(\boldsymbol{\theta})$ around the current iterate $\boldsymbol{\theta}_{t}$ as
\begin{equation}
\bar{J}_{i}^{t}\left(\boldsymbol{\theta}\right)=\hat{J}_{i}^{t}+\left(\hat{\boldsymbol{g}}_{i}^{t}\right)^{\intercal}\left(\boldsymbol{\theta}-\boldsymbol{\theta}_{t}\right)+\zeta_{i}\left\Vert \boldsymbol{\theta}-\boldsymbol{\theta}_{t}\right\Vert _{2}^{2},\forall i,\label{eq:surrogate functions}
\end{equation}
where $\zeta_{i}>0$ is a fixed regularization parameter that ensures strong convexity. The terms $\hat{J}_{i}^{t}$ and $\hat{\boldsymbol{g}}_{i}^{t}$ are smoothed estimates of the function value $J_{i}(\boldsymbol{\theta}_{t})$ and its gradient $\nabla J_{i}(\boldsymbol{\theta}_{t})$, respectively. They are obtained by recursively averaging the instantaneous noisy estimates $\tilde{J}_{i}^{t}$ (from (\ref{eq:J-new})) and $\tilde{\boldsymbol{g}}_{i}^{t}$ (from (\ref{eq:gtlt})) as
\begin{equation}
\hat{J}_{i}^{t}=\left(1-\alpha_{t}\right)\hat{J}_{i}^{t-1}+\alpha_{t}\tilde{J}_{i}^{t},\forall i,\label{eq:J-average}
\end{equation}

\begin{equation}
\hat{\boldsymbol{g}}_{i}^{t}=\left(1-\alpha_{t}\right)\hat{\boldsymbol{g}}_{i}^{t-1}+\alpha_{t}\tilde{\boldsymbol{g}}_{i}^{t},\forall i,\label{eq:g-average}
\end{equation}
where $\{\alpha_{t}\}$ is a decreasing step-size sequence satisfying Assumption 2 in Section \ref{sec:CONVERGENCE-ANALYSIS}. This iterative averaging is crucial as it reduces the variance of the stochastic estimates, leading to more stable surrogate functions. Combined with the mixed offline-online data reuse in (\ref{eq:g-tilde-1}), (\ref{eq:g-tilde}) and (\ref{eq:J-new}), it forms a dual mechanism for variance reduction.

Step 2: Solving the Surrogate Optimization Problem Using the convex surrogates $\{\bar{J}_{i}^{t}(\boldsymbol{\theta})\}$, we form a surrogate optimization problem. The primary goal is to minimize the surrogate objective while satisfying the surrogate constraints:
\begin{align}
\bar{\boldsymbol{\theta}}_{t}=\underset{\boldsymbol{\theta}\in\mathbf{\Theta}}{\textrm{argmin}}\  & \bar{J}_{0}^{t}\left(\boldsymbol{\theta}\right)\label{eq:objective update}\\
\textrm{s.t.}\  & \bar{J}_{i}^{t}\left(\boldsymbol{\theta}\right)\leq0,i=1,\cdots,I.\nonumber
\end{align}
In early iterations or with poor initial estimates, problem (\ref{eq:objective update}) might be infeasible. To guarantee progress, we employ a feasibility restoration step. We introduce an auxiliary slack variable $y$ and solve:
\begin{align}
\bar{\boldsymbol{\theta}}_{t}=\underset{\boldsymbol{\theta}\in\mathbf{\Theta},y}{\textrm{argmin}}\  & y\label{eq:feasible update}\\
\textrm{s.t.}\  & \bar{J}_{i}^{t}\left(\boldsymbol{\theta}\right)\leq y,i=1,\cdots,I.\nonumber
\end{align}
This problem always has a feasible solution and minimizes the maximum constraint violation. When (\ref{eq:objective update}) is feasible, the solution to (\ref{eq:feasible update}) naturally yields $y\leq0$ and coincides with the feasible solution.

Both (\ref{eq:objective update}) and (\ref{eq:feasible update}) are convex quadratic programs (QPs) with a strongly convex objective due to the $\zeta_{i}\|\cdot\|^{2}$ term. This structure guarantees a unique solution and permits efficient solving via standard convex optimization techniques, such as Lagrange-dual methods, often with closed-form updates.

Step 3: Policy Parameter Update Finally, the policy parameter is updated by taking a step from the current iterate $\boldsymbol{\theta}_{t}$ towards the solution of the surrogate problem $\bar{\boldsymbol{\theta}}_{t}$:
\begin{equation}
\boldsymbol{\theta}_{t+1}=(1-\beta_{t})\boldsymbol{\theta}_{t}+\beta_{t}\bar{\boldsymbol{\theta}}_{t}.\label{eq:theta update}
\end{equation}
where $\{\beta_{t}\}$ is a decreasing step-size sequence satisfying Assumption 2 in Section \ref{sec:CONVERGENCE-ANALYSIS}. This update is a convex combination, ensuring that $\boldsymbol{\theta}_{t+1}$ remains within the feasible parameter set $\boldsymbol{\Theta}$ if $\boldsymbol{\Theta}$ is convex. The sequence $\{\bar{\boldsymbol{\theta}}_{t}\}$ acts as a guide, and $\{\beta_{t}\}$ controls the learning rate, ensuring stable convergence towards a Karush-Kuhn-Tucker (KKT) point of the original problem, as will be proven in Section \ref{sec:CONVERGENCE-ANALYSIS}.

Summary: The actor module transforms a hard, stochastic, non-convex constrained problem into a sequence of easy-to-solve convex QPs. The CSSCA framework, powered by variance-reduced gradient estimates and safeguarded by feasibility restoration, provides a robust and theoretically sound mechanism for policy improvement under constraints.
\begin{algorithm}
\caption{Fused-CPRO Algorithm}\label{alg:the-Single-Loop-Deep}
\begin{algorithmic}[1]
\STATE \textbf{Input:} decreasing sequences $\{\xi_{t}\}$, $\{\alpha_{t}\}$, $\{\beta_{t}\}$, $\{\eta_{t}\}$, and $\{\gamma_{t}\}$.
\STATE Randomly initialize $\boldsymbol{\omega}_{0}^{i}$ and $\boldsymbol{\theta}_{0}$ from $\mathcal{N}(0,1/m^{2})$.
\FOR{$t=0,1,\cdots$}
\STATE Obtain $T_{\text{off}}^{t}$ offline samples $\mathcal{D}_{\text{off}}^{t}$ by sampling from $D_{1},\ldots,D_{N}$ according to $\boldsymbol{\rho}_{0}$.
\STATE Sample the new observation $\varepsilon_{t}$ and update the online dataset $\mathcal{D}_{0}^{t}$.
\STATE \textbf{Critic step:} update $\boldsymbol{\omega}_{t}^{i}$ and $\bar{\boldsymbol{\omega}}_{t}^{i}$ by (\ref{eq:TD-learning}) and (\ref{eq:w^_}), respectively.
\STATE \textbf{Actor step:} calculate $\hat{J}_{i}^{t}$ via (\ref{eq:J-new}) and (\ref{eq:J-average}).
\STATE Estimate $\hat{\boldsymbol{g}}_{i}^{t}$ via (\ref{eq:g-tilde}) and (\ref{eq:g-average}).
\STATE Update $\{\bar{J}_{i}^{t}(\boldsymbol{\theta})\}_{i=0,\ldots,I}$ via (\ref{eq:surrogate functions}).
\IF{Problem (\ref{eq:objective update}) is feasible}
\STATE Solve (\ref{eq:objective update}) to obtain $\bar{\boldsymbol{\theta}}_{t}$.
\ELSE
\STATE Solve (\ref{eq:feasible update}) to obtain $\bar{\boldsymbol{\theta}}_{t}$.
\ENDIF
\STATE Update $\boldsymbol{\theta}_{t+1}$ according to (\ref{eq:theta update}).
\ENDFOR
\end{algorithmic}
\end{algorithm}

\subsection{Comparative Analysis and Key Innovations}
Our proposed Fused-CPRO algorithm builds upon, yet significantly advances, the principles established in our prior convergent CRL works: the actor-only SCAOPO [citeSCAOPO] and the actor-critic SLDAC [xx]. While all three algorithms share a common foundation of convergence-provable design using CSSCA to handle non-convex constraints and employ outdated online data reuse for sample efficiency, Fused-CPRO introduces two fundamental innovations that define its superiority.

The primary evolutionary leap from SCAOPO to SLDAC was the adoption of an actor-critic architecture. Replacing the high-variance Monte Carlo value estimation with low-variance TD learning via DNNs provided more stable and accurate gradient estimates, which is a prerequisite for handling complex policy structures.

Fused-CPRO goes a decisive step further by introducing a principled policy reuse mechanism, which distinguishes it fundamentally from both predecessors. This leads to two concrete algorithmic differences from SLDAC:

\begin{enumerate}
\item Joint Optimization of Policy and Reuse Probabilities:

SLDAC optimizes only a single new policy. In contrast, Fused-CPRO jointly optimizes the parameters of the new DNN-based policy and the reuse probability vector $\boldsymbol{\rho}$ for all old policies (both TL and DK). This transforms policy selection from a heuristic into a learned component of the optimization, enabling dynamic, performance-driven knowledge fusion.
\item Systematic Offline Data Reuse:

Beyond reusing past online data, Fused-CPRO actively leverages offline datasets from all old policies. This allows the algorithm to extract gradient information from historical interactions in similar environments before costly online exploration begins, drastically reducing the initial interaction burden.
\end{enumerate}

These innovations directly translate into the key advantages demonstrated in the simulations: faster convergence (due to warm-starting from offline data and adaptive policy selection) and significantly lower online interaction cost. However, they also introduce considerable analytical complexity. The interdependence between the learning policy and the reuse probabilities, coupled with the potential bias introduced by offline data distribution shift, makes the convergence and performance guarantee analysis for Fused-CPRO markedly more challenging than for SLDAC. Overcoming these theoretical hurdles, as accomplished in Section \ref{sec:CONVERGENCE-ANALYSIS}, constitutes a major contribution of this work.

In summary, Fused-CPRO is not merely an incremental improvement but a substantial architectural advance that unifies heterogeneous knowledge fusion, constrained optimization, and sample-efficient learning within a single, provably convergent framework.

\section{Convergence and Performance Analysis}
\label{sec:CONVERGENCE-ANALYSIS}
In the following, we begin by stating the key assumptions for convergence analysis and then present the convergence rate of the critic module. Based on this, we further show that the surrogate functions satisfy asymptotic consistency. Finally, we prove that the proposed Algorithm \ref{alg:the-Single-Loop-Deep} converges to a KKT point.

\subsection{Key Assumptions and New Technical Challenges}
The convergence analysis employs a similar approach as that used in the single-loop deep actor-critic algorithm with CSSCA. Specifically, we first demonstrate that the critic accurately tracks the surrogate Q-functions with diminishing error. We then establish asymptotic consistency for the surrogate function values and corresponding surrogate gradients utilized in the actor update. Finally, the classical CSSCA framework is applied to conclude the convergence to a KKT point. Compared to the conventional single-policy method, the proposed Fused-CPRO algorithm introduces two additional bias terms, posing new challenges that require explicit control:

\begin{itemize}
\item \textbf{Policy reuse optimization:}

The mixed policy $\pi_{\boldsymbol{\theta}}(\boldsymbol{a}\mid\boldsymbol{s})=\sum_{n=0}^{N}\rho_{n}\pi_{n}(\boldsymbol{a}\mid\boldsymbol{s})$ induces policy-gradient terms that involve mixture ratios of the form $\pi_{n}(\boldsymbol{a}\mid\boldsymbol{s})/\pi_{\boldsymbol{\theta}}(\boldsymbol{a}\mid\boldsymbol{s})$. To ensure numerical stability and avoid ill-conditioned gradient estimates, these ratios are required to be uniformly bounded..
\item \textbf{Offline dataset reuse:}

The mixed-data estimators introduced in equations \eqref{eq:g-tilde-1}--\eqref{eq:J-new} combine online samples with offline samples collected from distributions different from that of the current policy. This leads to a persistent distribution mismatch. The proposed Fused-CPRO employs a diminishing offline weight sequence $\{\xi_{t}\}$ to mitigate this mismatch.
\end{itemize}
We next present the key assumptions required for the convergence analysis for Fused-CPRO algorithm. These assumptions are standard in the analysis of infinite horizon average-cost methods with deep actor-critic framework.

\begin{assumption}[Assumptions on the Problem Structure]~

\label{assumption:problem_structure}

\begin{itemize}
\item \textbf{Ergodicity and geometric mixing:}

For any feasible parameter $\boldsymbol{\theta}\in\boldsymbol{\Theta}$,the Markov chain induced by $\pi_{\boldsymbol{\theta}}$ admits a unique stationary state-action distribution $\mathbf{P}_{\pi_{\boldsymbol{\theta}}}$. For all $t=0,1,...,$ the total-variation distance $d_{\mathrm{TV}}$ between $\mathbf{P}^{t}(\cdot\mid\boldsymbol{s})$ and $\mathbf{\mathbf{P}_{\pi_{\boldsymbol{\theta}}}}$ has upper bound: there exist constants $\lambda>0$ and $\varrho\in(0,1)$ such that
\[
\sup_{\boldsymbol{s}\in\mathcal{S}}d_{\mathrm{TV}}\!\left(\mathbf{P}^{t}(\cdot\mid\boldsymbol{s}),\,\mathbf{P}_{\pi_{\boldsymbol{\theta}}}\right)\leq\lambda\varrho^{t},\ \forall t\geq0.
\]
\item \textbf{Compactness and boundedness:}

The state space $\mathcal{S}$ and the action space $\mathcal{A}$ are both compact. And the costs/rewards $\{C_{i}^{\text{'}}(\boldsymbol{s},\boldsymbol{a})\}$ are uniformly bounded. The DNN's parameter space $\boldsymbol{\Theta}$ and $\Omega_{i}$ are compact and convex. The output of all DNNs are uniformly bounded. Also, the policy $\pi_{\theta}$ follows Lipschitz continuity over parameter $\boldsymbol{\theta}\in\boldsymbol{\Theta}$. 

\item \textbf{Offline data admissibility:}

Each offline dataset $\mathcal{D}_{n}$ is generated under the same CMDP dynamics and cost functions as the target task, but by an old policy $\pi_{n}$. The resulting distribution mismatch is bounded in the sense that the induced stationary distributions are within a finite total-variation radius.
\end{itemize}
\end{assumption}

\begin{assumption}[Step sizes and offline weight schedule]~

\label{assumption:step_size}

The sequences $\{\eta_{t}\}$, $\{\gamma_{t}\}$, $\{\alpha_{t}\}$, $\{\beta_{t}\}$, and $\{\xi_{t}\}$ are deterministic, positive, and non-increasing.

\begin{enumerate}
\item The sequence $\{\alpha_{t}\}$ satisfies $\alpha_{t}\to0$, $\sum_{t}\alpha_{t}=\infty$,and $\sum_{t}\alpha_{t}^{2}<\infty$.
\item The sequence $\{\beta_{t}\}$ satisfies $\beta_{t}\to0$, $\sum_{t}\beta_{t}=\infty$,and $\sum_{t}\beta_{t}^{2}<\infty$. In addition, the actor update runs on a slower time-scale, i.e., $\lim_{t\to\infty}\beta_{t}/\alpha_{t}=0$, and the window-drift term is summable: $\sum_{t=0}^{\infty}\alpha_{t}\beta_{t}\log(t+1)<\infty$.
\item The critic step size $\{\eta_{t}\}$ and the target-network step size $\{\gamma_{t}\}$ satisfy $\eta_{t}\to0,\gamma_{t}\to0$ and $\sum_{t}\eta_{t}=\infty,\sum_{t}\gamma_{t}=\infty$.
\item The actor evolves on the slowest time scale:
\[
\lim_{t\rightarrow\infty}\frac{\beta_{t}}{\alpha_{t}}=0,\qquad\lim_{t\rightarrow\infty}\frac{\beta_{t}}{\eta_{t}}=0.
\]
\item The offline weight $\xi_{t}\to0$ and satisfies an extra condition with the tracking step size, i.e., $\sum_{t}\alpha_{t}\xi_{t}<\infty$, so that the cumulative bias introduced by offline reuse is finite.
\item Define $n_{t}=t-\lceil t^{0.43}\rceil$.Then the following series are finite:

\begin{flalign}
& \sum_{t=0}^{\infty}\alpha_{t}(1-\gamma_{t})^{t^{0.215}}\gamma_{n_{t}}^{-1/2}<\infty,\sum_{t=0}^{\infty}\alpha_{t}\gamma_{n_{t}}^{1/2}\eta_{n_{t}}^{-1/2}<\infty, &&\nonumber \\
& \sum_{t=0}^{\infty}m_{Q}\alpha_{t}\gamma_{n_{t}}^{1/2}\eta_{n_{t}}^{1/2}t^{0.215}<\infty,\sum_{t=0}^{\infty}m_{Q}\alpha_{t}\eta_{n_{t}}t^{-0.57}<\infty, &&\nonumber \\
& \sum_{t=0}^{\infty}m_{Q}\alpha_{t}\eta_{n_{t}}\beta_{n_{t}}t^{0.86}<\infty,\sum_{t=0}^{\infty}m_{Q}\alpha_{t}\gamma_{n_{t}}^{1/2}t^{0.215}\xi_{n_{t}}<\infty, &&
\end{flalign}
where $m_{Q}$ denotes the width of DNN used in critic module.
\end{enumerate}
\end{assumption}

It is worth noting that Assumption \ref{assumption:step_size}.1- \ref{assumption:step_size}.3 is mentioned in our previous works of SLDAC. The additional condition Assumption \ref{assumption:step_size}.5 is specific to mixed-data reuse, ensuring that the distribution-mismatch bias contributed by offline samples is summable and does not accumulate over time. If we set the step size as following order:

\begin{align}
\alpha_{t} & =\mathcal{O}\!\big(m_{Q}^{-1/2}(t+1)^{-\kappa_{1}}\big), & \beta & =\mathcal{O}\!\big(m_{Q}^{-1/2}(t+1)^{-\kappa_{2}}\big),\nonumber \\
\eta_{t} & =\mathcal{O}\!\big(m_{Q}^{-1/2}(t+1)^{-\kappa_{3}}\big), & \gamma_{t} & =\mathcal{O}\!\big((t+1)^{-\kappa_{4}}\big),\nonumber \\
\xi_{t} & =\mathcal{O}\!\big((t+1)^{-\kappa_{5}}\big),\label{eq:poly_stepsize}
\end{align}

then Assumption \ref{assumption:step_size} can be satisfied when $\kappa_{1},\kappa_{2},\kappa_{3},\kappa_{4},\kappa_{5}$ lie in the following region:

\begin{equation}
\left\{ \begin{aligned} & 1>2\kappa_{2}-1>\kappa_{1}>0.43>\kappa_{4}>0,\\
& \min\{0.5\kappa_{1}+0.5\kappa_{2}+0.5\kappa_{3}-0.5,\ \kappa_{1}+\kappa_{3}\}>0.43,\\ & \kappa_{1}+0.5\kappa_{4}-0.5\kappa_{3}>1,\\ & \kappa_{1}+0.5\kappa_{3}+0.5\kappa_{4}>1.215,\\ & \kappa_{1}+\kappa_{5}>1,\\ & \kappa_{3}+2\kappa_{5}>1,\\ & \kappa_{1}+\kappa_{5}+0.5\kappa_{5}>1.215.
\end{aligned}
\right.\label{eq:kappa_region_core}
\end{equation}

One typical choice of step size is given below. Choose $\kappa_{1}\in(0.5,1)$ and then pick $\kappa_{2}\in(0.5+0.5\kappa_{1},1)$. Next, select $\kappa_{3}\in(\max\{0.43-\kappa_{1},\ 1.92-\kappa_{1}-\kappa_{2}\},\ 1)$ and $\kappa_{4}\in(\max\{2+\kappa_{3}-2\kappa_{1},\ 2.43-2\kappa_{1}-\kappa_{3}\},\ 0.43)$. Finally, choose $\kappa_{5}>\max\{1-\kappa_{1},\ (1-\kappa_{3})/2,\ 1.215-\kappa_{1}-0.5\kappa_{4}\}$. For example, $\left(\kappa_{1},\kappa_{2},\kappa_{3},\kappa_{4},\kappa_{5}\right)=(0.9,\,0.96,\,0.21,\,0.425,\,0.4)$ is a feasible point. In practice, these regions are sufficient (but not necessary) and can be mildly relaxed to improve transient learning speed, provided stability is maintained.

\subsection{Convergence of the Critic Module}
We establish a finite-time tracking bound for the critic module. The convergence analysis of the critic module differs from our prior work owing to the policy reuse and offline data. This module aims at evaluating the surrogate $Q$-functions $\{\hat{Q}_{i}^{\pi_{\boldsymbol{\theta}_{t}}}\}_{i=0}^{I}$ defined in $~\eqref{eq:definition of Qhat}$ with DNNs $\{f(\bar{\boldsymbol{\omega}}^{i};\boldsymbol{s},\boldsymbol{a})\}_{i=0}^{I}$ under the current mixed policy $\pi_{\boldsymbol{\theta}_{t}}$. Also, the critic module uses the running estimate $\hat{J}_{i}^{t}$ in place of the unknown average reward/cost. The estimated error can be derived as:
\begin{equation}
\epsilon_{\mathrm{cri}}(t)=\bigl|\mathbb{E_{\mathrm{\mathit{p_{t}}}}}[f(\bar{\boldsymbol{\omega}}_{t}^{i})|\bar{\boldsymbol{\omega}}_{t-1}^{i}]-\hat{Q}_{i}^{\pi_{\boldsymbol{\theta}_{t}}}(\boldsymbol{s},\boldsymbol{a})\bigr|,\forall i\label{eq:critic_error}
\end{equation}
For each critic DNN with with initialization $\boldsymbol{\omega}_{0}^{i}$ and constraint set $\Omega_{i}\triangleq\mathbb{B}(\boldsymbol{\omega}_{0}^{i},R_{\boldsymbol{\omega}})$, its local linearization function class can be defined as:
\begin{equation}
\hat{f}(\boldsymbol{\omega}^{i};\boldsymbol{s},\boldsymbol{a})\triangleq f(\boldsymbol{\omega}_{0}^{i};\boldsymbol{s},\boldsymbol{a})+\bigl\langle\nabla_{\boldsymbol{\omega}}f(\boldsymbol{\omega}_{0}^{i};\boldsymbol{s},\boldsymbol{a}),\,\boldsymbol{\omega}^{i}-\boldsymbol{\omega}_{0}^{i}\bigr\rangle,\quad\boldsymbol{\omega}^{i}\in\Omega_{i},\label{eq:local_linearization}
\end{equation}
and the induced function class is defined as $\hat{\mathcal{F}}_{i}\triangleq\{\hat{f}(\boldsymbol{\omega}^{i};\cdot):\boldsymbol{\omega}^{i}\in\Omega_{i}\}$. This linearization satisfies the identity
\begin{equation}
\hat{f}(\boldsymbol{\omega}_{a}^{i};\boldsymbol{s},\boldsymbol{a})-\hat{f}(\boldsymbol{\omega}_{b}^{i};\boldsymbol{s},\boldsymbol{a})=\bigl\langle\nabla_{\boldsymbol{\omega}}f(\boldsymbol{\omega}_{0}^{i};\boldsymbol{s},\boldsymbol{a}),\,\boldsymbol{\omega}_{a}^{i}-\boldsymbol{\omega}_{b}^{i}\bigr\rangle,\quad\forall i.
\end{equation}

\begin{assumption}[Assumptions on the target Q]~
\label{assumption:target_Q}
\begin{enumerate}
\item \textbf{Represent ability:}

For each $t$ and $i$,there exists a point $\dot{\boldsymbol{\omega}}_{t}^{i}\in\Omega_{i}$ such that $\hat{f}(\dot{\boldsymbol{\omega}}_{t}^{i};\boldsymbol{s},\boldsymbol{a})=\hat{Q}_{i}^{\pi_{\boldsymbol{\theta}_{t}}}(\boldsymbol{s},\boldsymbol{a})$.
\item \textbf{Uniform TD stability:}

The TD feature matrix is defined as:
\begin{align}
\mathbf{A}_{\boldsymbol{\theta}_{t}} & \triangleq\mathbb{E}_{p_{t}}\Bigl[\nabla_{\boldsymbol{\omega}}f(\boldsymbol{\omega}_{0}^{i};\boldsymbol{s}_{t},\boldsymbol{a}_{t})\nonumber \\
& \bigl(\nabla_{\boldsymbol{\omega}}f(\boldsymbol{\omega}_{0}^{i};\boldsymbol{s}_{t},\boldsymbol{a}_{t})-\nabla_{\boldsymbol{\omega}}f(\boldsymbol{\omega}_{0}^{i};\boldsymbol{s}_{t+1},\boldsymbol{a}_{t+1}')\bigr)^{\intercal}\Bigr]
\end{align}
where $\boldsymbol{a}_{t+1}'\sim\pi_{\boldsymbol{\theta}_{t}}(\cdot\mid\boldsymbol{s}_{t+1})$. The matrix $\mathbf{A}_{\boldsymbol{\theta}_{t}}$ satisfies $\mathbf{\boldsymbol{\omega}^{\mathit{T}}A}_{\boldsymbol{\theta}_{t}}\boldsymbol{\omega}>0$ for all $\boldsymbol{\omega}$. Also, the symmetric part $\mathbf{A}_{\boldsymbol{\theta}_{t}}+\mathbf{A}_{\boldsymbol{\theta}_{t}}^{\intercal}$ is uniformly positive definite over $\boldsymbol{\theta}_{t}\in\boldsymbol{\Theta}$, i.e., $\lambda_{min}(\mathbf{A}_{\boldsymbol{\theta}_{t}}+\mathbf{A}_{\boldsymbol{\theta}_{t}}^{\intercal})>\varsigma$ for constant $\varsigma>0$.
\end{enumerate}
\end{assumption}

Based on the mentioned assumptions and definitions, we can obtain the convergence rate of the critic module by decomposing its tracking error. Firstly, \eqref{eq:critic_error} can be decomposed and upper bounded by the sum of the linearization error and the approximation error between the linearized function $\hat{f}(\bar{\boldsymbol{\omega}}_{t}^{i})$ and the true Q-function $\hat{Q}_{i}^{\pi_{\theta_{t}}}$:

\begin{equation}
\epsilon_{t,i}^{\mathrm{cri}}\le\left|\mathbb{E}\big[f(\bar{\boldsymbol{\omega}}_{t}^{i})\big]-\mathbb{E}\big[\hat{f}(\bar{\boldsymbol{\omega}}_{t}^{i})\big]\right|+\left|\mathbb{E}\big[\hat{f}(\bar{\boldsymbol{\omega}}_{t}^{i})\big]-\hat{Q}_{i}^{\pi_{\theta_{t}}}\right|,
\end{equation}
where the first term will be bounded in \ref{sec:Technical-Bounds-DNN} and the second term can be further expressed as follows:

\begin{equation}
\left|\mathbb{E}\big[\hat{f}(\bar{\boldsymbol{\omega}}_{t}^{i})\big]-\hat{Q}_{i}^{\pi_{\theta_{t}}}\right|\le\|\nabla_{\omega}f(\boldsymbol{\omega}_{0})\|_{2}\cdot\left\Vert \mathbb{E}[\bar{\boldsymbol{\omega}}_{t}^{i}]-\dot{\omega}_{t}^{i}\right\Vert _{2},\ \forall i,
\end{equation}
where $\|\nabla_{\omega}f(\boldsymbol{\omega}_{0})\|_{2}$ will be bounded in Appendix \ref{sec:Technical-Bounds-DNN}. To decouple the error, we use a time separation method to form a surrogate trajectory. We introduce two auxiliary parameters $\{\bar{\boldsymbol{m}}_{t}^{i},\boldsymbol{m}_{t}^{i},\forall i\}$:
\begin{equation}
\boldsymbol{m}_{k}^{i}=\Pi_{\Omega_{i}}\!\left(\boldsymbol{m}_{k-1}^{i}-\eta_{k}\boldsymbol{M}_{k-1}^{i}\right),\bar{\boldsymbol{m}}_{k}^{i}=(1-\gamma_{k})\bar{\boldsymbol{m}}_{k-1}^{i}+\gamma_{k}\boldsymbol{m}_{k}^{i},
\end{equation}
where $\boldsymbol{M}_{k-1}^{i}$ is formulated as:
\begin{align}
\boldsymbol{M}_{k-1}^{i} & =\Bigl(\hat{f}(\boldsymbol{m}_{k-1}^{i};\tilde{\boldsymbol{s}}_{k},\tilde{\boldsymbol{a}}_{k})-\bigl(C_{i}^{\text{'}}(\tilde{\boldsymbol{s}}_{k},\tilde{\boldsymbol{a}}_{k})-\hat{J}_{i}^{k-1}\nonumber \\
& +\hat{f}(\boldsymbol{m}_{k-1}^{i};\tilde{\boldsymbol{s}}_{k+1},\tilde{\boldsymbol{a}}_{k+1}')\bigr)\Bigr)\nabla_{\boldsymbol{\omega}}f(\boldsymbol{\omega}_{0}^{i};\tilde{\boldsymbol{s}}_{k},\tilde{\boldsymbol{a}}_{k}),
\end{align}
Along the surrogate trajectory, the timeline is partitioned into two segments by the index $n_{t}\triangleq t-\lceil t^{\kappa_{6}}\rceil$. The surrogate observations $\{\tilde{\varepsilon}_{k}^{(t)}\}$ can be presented as:
\begin{equation}
\tilde{\varepsilon}_{k}^{(t)}=\bigl(\tilde{\boldsymbol{s}}_{k},\tilde{\boldsymbol{a}}_{k},\{C_{i}^{\text{'}}(\tilde{\boldsymbol{s}}_{k},\tilde{\boldsymbol{a}}_{k})\}_{i=0}^{I},\tilde{\boldsymbol{s}}_{k+1}\bigr).
\end{equation}
The observations after $n_{t}$ are sampled with fixed policy $\pi_{\boldsymbol{\theta}_{n_{t}}}$ at $t=n_{t}$, while the observations before $n_{t}$ are sampled with varying policy $\pi_{\boldsymbol{\theta}_{k}}(\cdot\mid\tilde{\boldsymbol{s}}_{k})$ for each time step $k$. Then, $\left\Vert \mathbb{E}[\bar{\boldsymbol{\omega}}_{t}^{i}]-\dot{\boldsymbol{\omega}}_{t}^{i}\right\Vert _{2}$ can be bounded as the summation of two parts:
\begin{equation}
\left\Vert \mathbb{E}[\bar{\boldsymbol{\omega}}_{t}^{i}]-\dot{\boldsymbol{\omega}}_{t}^{i}\right\Vert _{2}\le\left\Vert \mathbb{E}[\bar{\boldsymbol{\omega}}_{t}^{i}]-\mathbb{E}[\bar{\boldsymbol{m}}_{t}^{i}]\right\Vert _{2}+\left\Vert \mathbb{E}[\bar{\boldsymbol{m}}_{t}^{i}]-\dot{\boldsymbol{\omega}}_{t}^{i}\right\Vert _{2},
\end{equation}
where the first term is the error of the finite-time TD tracking on a frozen window and the second term is the mismatch due to policy drift. 

\begin{lemma}[Convergence Rate of critic]~

\label{lemma-convergence-rate-critic}

Fix $\kappa_{6}\in(0,1)$ and define $n_{t}\triangleq t-\lceil t^{\kappa_{6}}\rceil$. Under the assumptions stated above, the critic error admits the bound
\begin{align}
\epsilon_{t,i}^{\mathrm{cri}} & \le\mathcal{O}\Bigl(\epsilon_{m_{Q}}+\frac{(1-\gamma_{t})^{t^{\kappa_{6}/2}}}{\sqrt{\gamma_{t}}}+m_{Q}\sqrt{\gamma_{n_{t}}\eta_{n_{t}}}t^{\kappa_{6}/2}+\nonumber \\
& +m_{Q}\eta_{n_{t}}t^{\kappa_{6}-1}+\sqrt{\tfrac{\gamma_{n_{t}}}{\eta_{n_{t}}}}\nonumber \\ & m_{Q}\eta_{n_{t}}\beta_{n_{t}}t^{2\kappa_{6}}+m_{Q}\sqrt{\gamma_{n_{t}}}t^{\kappa_{6}/2}\,\xi_{n_{t}}\Bigr),\quad\forall i,
\end{align}
with almost probability 1. Here $\xi_{t}$ is the offline weight used in the mixed estimation of $\hat{J}_{i}^{t}$. $\epsilon_{m_{Q}}$ is the local linearization error. Consequently, if $\xi_{t}\to0$ and the step sizes are chosen so that every term on the right-hand side vanishes, then $\epsilon_{t,i}^{\mathrm{cri}}\to\mathcal{O}(\epsilon_{m_{Q}})$. Moreover, $\epsilon_{m_{Q}}\to0$ as $m_{Q}\to\infty$.

\end{lemma}

\subsection{Asymptotic Consistency of the Surrogate Functions}
The actor update relies on the running estimates $\hat{J}_{i}^{t}$ and $\hat{\boldsymbol{g}}_{i}^{t}$ in (\ref{eq:J-average})--(\ref{eq:g-average}). Observation/data reuse in the online buffer induces an off-policy bias because samples are drawn under slightly outdated policy parameters. Also, the offline data used in Fused-CPRO introduces offline bias. Compared with prior work in SLDAC and SACOPO, the offline bias is more challenging: its distribution mismatch does not vanish by itself and must be suppressed by the decay of $\xi_{t}$. Despite the above challenges, we prove that the running estimates $\hat{J}_{i}^{t}$ and $\hat{\boldsymbol{g}}_{i}^{t}$ satisfy asymptotic consistency, which is presented as the lemma below.

\begin{lemma}[Asymptotic consistency of running estimates]~

\label{lemma-asymptotic-consistency}

Suppose Assumptions \ref{assumption:problem_structure}-\ref{assumption:step_size} hold. Then, for each $i\in\{0,1,\ldots,I\}$, the running estimates satisfy
\begin{align}
\lim_{t\rightarrow\infty}\bigl|\hat{J}_{i}^{t}-J_{i}(\boldsymbol{\theta}_{t})\bigr| & =0,\\
\lim_{t\rightarrow\infty}\bigl\|\hat{\boldsymbol{g}}_{i}^{t}-\nabla_{\boldsymbol{\theta}}J_{i}(\boldsymbol{\theta}_{t})\bigr\|_{2} & \leq c_{\pi}\epsilon_{m_{Q}},
\end{align}
almost surely, where $c_{\pi}>0$ is a constant independent of $t$ and $m_{Q}$.

\end{lemma}

\begin{IEEEproof}
The recursions in (\ref{eq:J-average}) and (\ref{eq:g-average}) fall into projected stochastic approximation with diminishing step size $\alpha_{t}$. Asymptotic consistency of $\hat{J}_{i}^{t}$ and $\hat{\boldsymbol{g}}_{i}^{t}$ can be established by verifying that the required technical conditions are satisfied. A complete proof is provided in Appendix \ref{sec:proof-of-asymptotic-consistency}.

\end{IEEEproof}

\subsection{Finite-Time Convergence Rate of the Surrogate Functions}
Lemma \ref{lemma-asymptotic-consistency} shows that the running estimates $\hat{J}_{i}^{t}$ and $\hat{\boldsymbol{g}}_{i}^{t}$ are asymptotically accurate during iteration. In this subsection, we provide a finite-time characterization that quantifies the convergence rate of the surrogate functions to make the roles of the buffer length $T_{t}$ and the offline mixing weight $\xi_{t}$ in the proposed Fused-CPRO algorithm more explicit.

We measure the estimation quality via the finite-time average estimation error:
\begin{align}
\epsilon_{J}(t) & \triangleq\frac{1}{t+1}\sum_{k=0}^{t}\max_{0\le i\le I}\mathbb{E}\!\left[\left|\hat{J}_{i}^{k}-J_{i}(\boldsymbol{\theta}_{k})\right|^{2}\right],\label{eq:epsJ_def}\\
\epsilon_{g}(t) & \triangleq\frac{1}{t+1}\sum_{k=0}^{t}\max_{0\le i\le I}\mathbb{E}\!\left[\left\Vert \hat{\boldsymbol{g}}_{i}^{k}-\nabla_{\boldsymbol{\theta}}J_{i}(\boldsymbol{\theta}_{k})\right\Vert _{2}^{2}\right].\label{eq:epsg_def}
\end{align}

\begin{lemma}[Convergence rate of the Surrogate Functions]~

\label{lemma-convergence-rate-surrogate}

Suppose Assumptions \ref{assumption:problem_structure}- \ref{assumption:target_Q} hold and the online buffer length $T_{t}=\mathcal{O}(\log(t+1))$. The finite-time convergence rate of the surrogate functions can be formulated as:

\begin{align}
\epsilon_{J}(t)\le & \ \mathcal{O}\!\Big(\frac{1}{(t+1)\alpha_{t+1}}\Big)+\mathcal{O}\!\Big(\frac{1}{t+1}\sum_{k=0}^{t}\big(\alpha_{k+1}+T_{k}\alpha_{k}\nonumber \\
& +T_{k}^{2}\beta_{k}+\beta_{k}^{2}(\alpha_{k+1}^{-2}+\alpha_{k+1}^{-1})\big)\Big)\nonumber \\ & +\mathcal{O}\!\Big(\frac{1}{t+1}\sum_{k=0}^{t}\xi_{k+1}\Big),\label{eq:epsJ_rate}\\[1mm]
\epsilon_{g}(t)\le & \ \mathcal{O}\!\Big(\frac{1}{(t+1)\alpha_{t+1}}\Big)+\mathcal{O}\!\Big(\frac{1}{t+1}\sum_{k=0}^{t}\big(\bar{\epsilon}_{\mathrm{cri}}(k)+k^{-1}\nonumber \\
& +T_{k}\beta_{k-T_{k}+1}+\beta_{k}^{2}(\alpha_{k+1}^{-2}+\alpha_{k+1}^{-1})+\alpha_{k+1}\big)\Big)\nonumber \\ & +\mathcal{O}\!\Big(\sqrt{\epsilon_{J}(t)}+\frac{1}{t+1}\sum_{k=0}^{t}\xi_{k+1}+\epsilon_{m_{Q}}\Big).
\end{align}
where the aggregated critic error $\bar{\epsilon}_{\mathrm{cri}}(k)\triangleq\max_{0\le i\le I}\epsilon_{k,i}^{\mathrm{cri}}$.

\end{lemma}

\begin{IEEEproof}
Both recursions in (\ref{eq:J-average}) and (\ref{eq:g-average}) can be viewed as stochastic approximation tracking the moving targets $J_{i}(\boldsymbol{\theta}_{t})$ and $\nabla_{\boldsymbol{\theta}}J_{i}(\boldsymbol{\theta}_{t})$. The analysis follows a decomposition of the error into:
\begin{itemize}
\item A martingale difference term induced by Markovian sampling.
\item A drift term due to the variation of $\boldsymbol{\theta}_{t}$ over the finite buffer window.
\item The critic tracking error.
\item The offline mismatch term weighted by $\xi_{t}$.
\end{itemize}
The detailed proof is provided in Appendix \ref{sec:proof-of-surrogate}.
\end{IEEEproof}
\subsection{Convergence of the Actor Module}
This section proves that the proposed Fused-CPRO algorithm converges to a KKT point of the original Problem \eqref{eq:problem 1-1}. Building upon the finite-time tracking bounds of the critic and the asymptotic consistency of the surrogate functions detailed in Lemma \ref{lemma-asymptotic-consistency}, we now analyze the limiting behavior of the policy parameter sequence $\{\boldsymbol{\theta}_{t}\}$. Firstly, we define the KKT solutions of the original problem \eqref{eq:problem 1-1}.

\begin{definition}
KKT solutions of the original problem.

A solution $\boldsymbol{\theta}^{\star}$ is called a KKT solution of the original problem, if there exist Lagrange multipliers $\lambda=\left[\lambda_{1},\ldots,\lambda_{I}\right]^{T}\succeq0$, such that the following conditions are satisfied:
\begin{gather}
g_{0}\!\left(\boldsymbol{\theta}^{\star}\right)+\sum_{i=1}^{I}\lambda_{i}g_{i}\!\left(\boldsymbol{\theta}^{\star}\right)=0,\label{eq:KKT_stationarity}\\
J_{i}\!\left(\boldsymbol{\theta}^{\star}\right)\leq0,\ i=1,\ldots,I,\label{eq:KKT_primal_feasibility}\\
\lambda_{i}J_{i}\!\left(\boldsymbol{\theta}^{\star}\right)=0,\ i=1,\ldots,I.\label{eq:KKT_complementary_slackness}
\end{gather}

\end{definition}
Moreover, considering a subsequence $\{\boldsymbol{\theta}_{t_{j}}\}_{j=1}^{\infty}$ converging to a limiting point $\boldsymbol{\theta}^{\star}$, there exist converged surrogate functions $\hat{J}_{i}(\boldsymbol{\theta}),\forall i$ such that
\begin{equation}
\lim_{j\rightarrow\infty}\bar{J}_{i}^{t_{j}}(\boldsymbol{\theta})=\hat{J}_{i}(\boldsymbol{\theta}),\ \forall\boldsymbol{\theta}\in\boldsymbol{\Theta},\label{eq:surrogate_limit_function}
\end{equation}
where
\begin{align}
\left|\hat{J}_{i}\!\left(\boldsymbol{\theta}^{\star}\right)-J_{i}\!\left(\boldsymbol{\theta}^{\star}\right)\right|&=0,\label{eq:surrogate_limit_value}\\
\left\Vert \nabla\hat{J}_{i}\!\left(\boldsymbol{\theta}^{\star}\right)-\nabla J_{i}\!\left(\boldsymbol{\theta}^{\star}\right)\right\Vert _{2}&=0.\label{eq:surrogate_limit_gradient}
\end{align}

Finally, with lemma \ref{lemma-convergence-rate-critic}-\ref{lemma-convergence-rate-surrogate} and assumptions mentioned, we can prove the global convergence theorem. 


\begin{theorem}[Global Convergence of Algorithm \ref{alg:the-Single-Loop-Deep}]~
\label{Theorem:GlobalConvergence}

Suppose Assumptions \ref{assumption:problem_structure}-\ref{assumption:target_Q} are satisfied and the initial point $\boldsymbol{\theta}_{0}$ is feasible, i.e., $\max_{i\in\{1,\ldots,I\}}J_{i}\!\left(\boldsymbol{\theta}_{0}\right)\leq0$, and the number of sampled data is set to $T_{t}=\mathcal{O}(\log t)$. Denote $\{\boldsymbol{\theta}_{t}\}_{t=1}^{\infty}$ as the iterates generated by Algorithm \ref{alg:the-Single-Loop-Deep} with a sufficiently small initial step size $\beta_{0}$. Then, every limiting point $\boldsymbol{\theta}^{\star}$ of $\{\boldsymbol{\theta}_{t}\}_{t=1}^{\infty}$ satisfying the Slater condition satisfies the KKT conditions in (\ref{eq:KKT_stationarity}), (\ref{eq:KKT_primal_feasibility}), and (\ref{eq:KKT_complementary_slackness}) up to an error of $\epsilon_{m_{Q}}$, i.e.,
\begin{gather}
\left\Vert g_{0}\!\left(\boldsymbol{\theta}^{\star}\right)+\sum_{i}\lambda_{i}g_{i}\!\left(\boldsymbol{\theta}^{\star}\right)\right\Vert _{2}\leq\epsilon_{m_{Q}},\label{eq:global_conv_stationarity_residual}\\
J_{i}\!\left(\boldsymbol{\theta}^{\star}\right)\leq\epsilon_{m_{Q}},\ i=1,\ldots,I,\label{eq:global_conv_primal_residual}\\
\lambda_{i}J_{i}\!\left(\boldsymbol{\theta}^{\star}\right)\leq\epsilon_{m_{Q}},\ i=1,\ldots,I.\label{eq:global_conv_complement_residual}
\end{gather}
where $\lim_{m_{Q}\rightarrow\infty}\epsilon_{m_{Q}}=0$.
\end{theorem}

We provide proof in appendix \ref{sec:proof-of-critic}-\ref{sec:proof-of-surrogate}. Leveraging the theoretical groundwork in our prior work, we detail only the parts of the proof that depart from existing analyses and are technically more involved, arising from the incorporation of offline data reuse and policy reuse in the proposed algorithm. 

\section{Simulation Results}
In this section, we provide numerical simulation results in three typical CMDP scenarios.  
\section{Conclusion}
In this paper, we introduced a novel CPRO framework to enhance the sample efficiency and robustness of constrained reinforcement learning. The core of this work is a unified formulation that integrates heterogeneous prior knowledge (encompassing both policies from transfer learning and those derived from domain expertise) into a single CMDP. This allows for the joint optimization of a new learnable policy and the adaptive reuse probabilities of all old policies. To solve this CMDP, we developed the Fused-CPRO algorithm, which employs a CSSCA-based actor for reliable optimization under non-convex stochastic constraints and a mixed offline-online data reuse critic for high sample efficiency. We provided rigorous theoretical analysis, proving that Fused-CPRO converges almost surely to a KKT point and delivers performance guarantees superior to all prior policies, ensuring robustness by leveraging interpretable domain knowledge as a safeguard. Simulation results across several benchmarks consistently demonstrate that Fused-CPRO achieves significantly faster convergence, lower online interaction costs, and superior final performance compared to state-of-the-art baselines. These findings validate the proposed framework as an effective and principled approach for developing data-efficient and reliable CRL agents in complex, safety-critical environments.
\appendices

\section{Technical Bounds about DNN}
\label{sec:Technical-Bounds-DNN}
In this section, we restate the DNN lemma used for our convergence analysis for the Fused-CPRO. 
\begin{lemma}[Technical Bounds about DNN]~
Let $\sigma\in(0,1)$, $d$, and $\{a_{i}\}_{i=0,1,\ldots}$ denote some universal constants that are independent of proposed algorithm's parameters. Then, for any $\sigma>0$, if the radius of the constraint set satisfies
\begin{align}
a_{1}d^{3/2}L^{-1}m_{Q}^{-3/4}\leq R_{\boldsymbol{\omega}}\leq a_{2}L^{-6}(\log m_{Q})^{-3},
\end{align}
and the width of the DNN satisfies
\begin{align}
m_{Q}\geq a_{3}\max\left\{ dL^{2}\log\left(\frac{m_{Q}}{\sigma}\right),\,R_{\boldsymbol{\omega}}^{-4/3}L^{-8/3}\log\left(\frac{m_{Q}}{R_{\boldsymbol{\omega}}\sigma}\right)\right\}, 
\end{align}
it holds that the bias between $f(\boldsymbol{\omega})$ and its local linearization $\hat{f}(\boldsymbol{\omega})$ satisfies
\begin{align}
\left|f(\boldsymbol{\omega};\boldsymbol{s},\boldsymbol{a})-\hat{f}(\boldsymbol{\omega};\boldsymbol{s},\boldsymbol{a})\right|
\leq a_{4}R_{\boldsymbol{\omega}}^{4/3}L^{4}\sqrt{m_{Q}\log m_{Q}}+a_{5}R_{\boldsymbol{\omega}}^{2}L^{5},
\end{align}
with probability at least $1-\sigma-exp\{-a_6 m_QR^{2/3}_\omega L\}$. 

Also, DNN's gradient is bounded as:
\begin{align}
\left\Vert \nabla_{\boldsymbol{\omega}}f(\boldsymbol{\omega};\boldsymbol{s},\boldsymbol{a})\right\Vert _{2}\leq a_{7}m_{Q}^{1/2},
\end{align}
with probability at least $1-L^{2}\exp\{-a_{8}m_{Q}R_{\boldsymbol{\omega}}^{2/3}L\}$. In the theoretical analysis, we assume all the critics have same initial parameters generated from $\mathcal{N}( 0,1/m^2_Q)$. 

\end{lemma}

\section{Proof of Lemma \ref{lemma-convergence-rate-critic}}
\label{sec:proof-of-critic}
In this section, we omit the superscript $i$ for simplicity since no confusion will arise. Recall the critic error definition in \eqref{eq:critic_error}:
\begin{equation}
\epsilon_{t}^{\mathrm{cri}}
\triangleq\bigl|\,\mathbb{E}_{p_t}\big[f(\bar{\boldsymbol{\omega}}_{t};\boldsymbol{s},\boldsymbol{a})\,\big|\,\bar{\boldsymbol{\omega}}_{t-1}\big]-\hat Q^{\pi_{\boldsymbol{\theta}_t}}(\boldsymbol{s},\boldsymbol{a})\,\bigr|.
\end{equation}
Using the local linearization \eqref{eq:local_linearization} and the triangle inequality,
\begin{align}
\epsilon_{t}^{\mathrm{cri}}
&\le \underbrace{\bigl|\,\mathbb{E}[f(\bar{\boldsymbol{\omega}}_{t})]-\mathbb{E}[\hat f(\bar{\boldsymbol{\omega}}_{t})]\,\bigr|}_{\mathrm{(bias\ 1)}}
+\bigl|\,\mathbb{E}[\hat f(\bar{\boldsymbol{\omega}}_{t})]-\hat Q^{\pi_{\boldsymbol{\theta}_t}}\,\bigr|.
\label{eq:appB_decomp1}
\end{align}
By the local linearization error bound provided in Appendix~\ref{sec:Technical-Bounds-DNN}, $\mathrm{bias\ 1}=\mathcal{O}(\epsilon_{m_Q})$.

According to Assumption~\ref{assumption:target_Q}.1, there exists a point $\dot{\boldsymbol{\omega}}_{t}\in\Omega$ such that
$\hat f(\dot{\boldsymbol{\omega}}_{t})\equiv \hat Q^{\pi_{\boldsymbol{\theta}_t}}$.
Thus the second term in \eqref{eq:appB_decomp1} becomes
\begin{equation}
\bigl|\,\mathbb{E}[\hat f(\bar{\boldsymbol{\omega}}_{t})]-\hat Q^{\pi_{\boldsymbol{\theta}_t}}\,\bigr|
=\bigl|\,\mathbb{E}[\hat f(\bar{\boldsymbol{\omega}}_{t})]-\hat f(\dot{\boldsymbol{\omega}}_{t})\,\bigr|.
\label{eq:appB_decomp2}
\end{equation}
Using the linear identity induced by \eqref{eq:local_linearization},
$\hat f(\boldsymbol{\omega}_a)-\hat f(\boldsymbol{\omega}_b)=\langle\nabla_{\boldsymbol{\omega}}f(\boldsymbol{\omega}_0;\cdot),\,\boldsymbol{\omega}_a-\boldsymbol{\omega}_b\rangle$,
we obtain
\begin{equation}
\bigl|\,\mathbb{E}[\hat f(\bar{\boldsymbol{\omega}}_{t})]-\hat f(\dot{\boldsymbol{\omega}}_{t})\,\bigr|
\le \bigl\|\nabla_{\boldsymbol{\omega}}f(\boldsymbol{\omega}_0)\bigr\|_2\cdot\bigl\|\mathbb{E}[\bar{\boldsymbol{\omega}}_{t}]-\dot{\boldsymbol{\omega}}_{t}\bigr\|_2.
\label{eq:appB_reduce_to_param}
\end{equation}
From Appendix~\ref{sec:Technical-Bounds-DNN}, $\|\nabla_{\boldsymbol{\omega}}f(\boldsymbol{\omega}_0)\|_2$ is bounded, then,
\begin{equation}
\epsilon_{t}^{\mathrm{cri}}
\le \mathcal{O}(\epsilon_{m_Q})+\mathcal{O}(\sqrt{m_Q})\cdot\bigl\|\mathbb{E}[\bar{\boldsymbol{\omega}}_{t}]-\dot{\boldsymbol{\omega}}_{t}\bigr\|_2.
\label{eq:appB_core_goal}
\end{equation}
It remains to upper bound $\|\mathbb{E}[\bar{\boldsymbol{\omega}}_{t}]-\dot{\boldsymbol{\omega}}_{t}\|_2$.

Following the similar trick in SLDAC, we fix $\kappa_6\in(0,1)$ and define $n_t\triangleq t-\lceil t^{\kappa_6}\rceil$ as in Lemma~\ref{lemma-convergence-rate-critic}.
Recall the auxiliary recursion (already defined before Lemma~\ref{lemma-convergence-rate-critic}):
\begin{equation}
\begin{aligned}
\boldsymbol{m}_{k}
&=\Pi_{\Omega}\!\left(\boldsymbol{m}_{k-1}-\eta_{k}\boldsymbol{M}_{k-1}\right),\\
\bar{\boldsymbol{m}}_{k}
&=(1-\gamma_{k})\bar{\boldsymbol{m}}_{k-1}+\gamma_{k}\boldsymbol{m}_{k},
\end{aligned}
\label{eq:appB_aux_update}
\end{equation}
where $\boldsymbol{M}_{k-1}$ is the linearized TD-gradient derived by the surrogate trajectory $\{\tilde\varepsilon_{k}^{(t)}\}$ mentioned before Lemma~\ref{lemma-convergence-rate-critic}. 
For $k\in\{n_t+1,\ldots,t\}$, the surrogate trajectory after $n_t$ is generated under the \emph{frozen} policy $\pi_{\boldsymbol{\theta}_{n_t}}$, while the real trajectory uses the slowly varying policy sequence $\{\pi_{\boldsymbol{\theta}_{k}}\}$.

By adding and subtracting $\mathbb{E}[\bar{\boldsymbol{m}}_{t}]$ and using $(a+b)^2\le 2a^2+2b^2$, $\|\mathbb{E}[\bar{\boldsymbol{\omega}}_{t}]-\dot{\boldsymbol{\omega}}_{t}\|_2$ can be decomposed as: 
\begin{equation}
\bigl\|\mathbb{E}[\bar{\boldsymbol{\omega}}_{t}]-\dot{\boldsymbol{\omega}}_{t}\bigr\|_2^2
\le 2\underbrace{\bigl\|\mathbb{E}[\bar{\boldsymbol{m}}_{t}]-\dot{\boldsymbol{\omega}}_{t}\bigr\|_2^2}_{\mathrm{(bias\ 2)}}
+2\underbrace{\bigl\|\mathbb{E}[\bar{\boldsymbol{\omega}}_{t}]-\mathbb{E}[\bar{\boldsymbol{m}}_{t}]\bigr\|_2^2}_{\mathrm{(bias\ 3)}}.
\label{eq:appB_bias23_split}
\end{equation}
In the following, we next bound bias~2 and bias~3.

\subsection{Bounds about bias~2}
Unfolding the recursion $\bar{\boldsymbol{m}}_{k}=(1-\gamma_{k})\bar{\boldsymbol{m}}_{k-1}+\gamma_{k}\boldsymbol{m}_{k}$, there exist weights
$w_{t,t'}\triangleq \gamma_{t'}\prod_{j=t'+1}^{t}(1-\gamma_j)$ for $t'=0,1,\ldots,t$ such that
\begin{equation}
\bar{\boldsymbol{m}}_{t}=\sum_{t'=0}^{t} w_{t,t'}\,\boldsymbol{m}_{t'},\qquad \sum_{t'=0}^{t} w_{t,t'}=1.
\label{eq:appB_bar_m_weights}
\end{equation}
With Jensen's inequality, bias~2 has upper bound: 
\begin{equation}
\bigl\|\mathbb{E}[\bar{\boldsymbol{m}}_{t}]-\dot{\boldsymbol{\omega}}_{t}\bigr\|_2^2
\le \sum_{t'=0}^{t} w_{t,t'}\,\bigl\|\mathbb{E}[\boldsymbol{m}_{t'}]-\dot{\boldsymbol{\omega}}_{t}\bigr\|_2^2.
\label{eq:appB_jensen_bias2}
\end{equation}
Split the sum into early part $t'\le n_t$ and recent part $t'\in\{n_t+1,\ldots,t\}$. Since $\boldsymbol{m}_{t'}\in\Omega=\mathbb{B}(\boldsymbol{\omega}_0,R_{\boldsymbol{\omega}})$ and $\dot{\boldsymbol{\omega}}_{t}\in\Omega$,
$\|\mathbb{E}[\boldsymbol{m}_{t'}]-\dot{\boldsymbol{\omega}}_{t}\|_2\le 2R_{\boldsymbol{\omega}}$. Moreover, using the monotonicity of $\{\gamma_t\}$,
\begin{equation}
\sum_{t'=0}^{n_t} w_{t,t'}
\le \prod_{j=n_t+1}^{t}(1-\gamma_j)
\le (1-\gamma_t)^{t-n_t}
\le (1-\gamma_t)^{t^{\kappa_6}}.
\label{eq:appB_weight_early}
\end{equation}
Therefore,
\begin{equation}
\sum_{t'=0}^{n_t} w_{t,t'}\,\bigl\|\mathbb{E}[\boldsymbol{m}_{t'}]-\dot{\boldsymbol{\omega}}_{t}\bigr\|_2^2
\le 4R_{\boldsymbol{\omega}}^{2}(1-\gamma_t)^{t^{\kappa_6}}
=\mathcal{O}\!\left(\frac{(1-\gamma_t)^{t^{\kappa_6}}}{m_Q\,\gamma_t}\right),
\label{eq:appB_bias2_early}
\end{equation}
where we used the standard NTK scaling $R_{\boldsymbol{\omega}}=\mathcal{O}(m_Q^{-1/2})$ in Appendix~\ref{sec:Technical-Bounds-DNN} and the fact that $\lim_{t\to\infty}\gamma_t\to 0$.

Define the tracking error during the frozen window as
$\boldsymbol{e}_{k}\triangleq \mathbb{E}[\boldsymbol{m}_{k}]-\dot{\boldsymbol{\omega}}_{t}$ for $k\in\{n_t+1,\ldots,t\}$.
The error can be expanded using the non-expansive property of projection oepration: 
\begin{equation}
\begin{aligned}
\|\boldsymbol{e}_{k+1}\|_2^2
&=\bigl\|\mathbb{E}[\boldsymbol{m}_{k+1}]-\dot{\boldsymbol{\omega}}_{t}\bigr\|_2^2\\
&\le \bigl\|\boldsymbol{e}_{k}-\eta_{k+1}\,\mathbb{E}[\boldsymbol{M}_{k}(\boldsymbol{m}_{k})]\bigr\|_2^2\\
&=\|\boldsymbol{e}_{k}\|_2^2-2\eta_{k+1}\big\langle \mathbb{E}[\boldsymbol{M}_{k}(\boldsymbol{m}_{k})],\,\boldsymbol{e}_{k}\big\rangle\\
&\quad+\eta_{k+1}^2\bigl\|\mathbb{E}[\boldsymbol{M}_{k}(\boldsymbol{m}_{k})]\bigr\|_2^2.
\end{aligned}
\label{eq:appB_recursion_raw}
\end{equation}
The term $\big\langle \mathbb{E}[\boldsymbol{M}_{k}(\boldsymbol{m}_{k})],\,\boldsymbol{e}_{k}\big\rangle$ can be presented as: 
\begin{equation}
\begin{aligned}
\big\langle \mathbb{E}[\boldsymbol{M}_{k}(\boldsymbol{m}_{k})],\,\boldsymbol{e}_{k}\big\rangle
&=\big\langle \mathbb{E}[\boldsymbol{M}_{k}(\boldsymbol{m}_{k})-\boldsymbol{M}_{k}(\dot{\boldsymbol{\omega}}_{t})],\,\boldsymbol{e}_{k}\big\rangle\\
&\quad+\big\langle \mathbb{E}[\boldsymbol{M}_{k}(\dot{\boldsymbol{\omega}}_{t})],\,\boldsymbol{e}_{k}\big\rangle.
\end{aligned}
\label{eq:appB_add_subtract}
\end{equation}
The first term is controlled by Assumption~\ref{assumption:target_Q}.2: 
there exists $\varsigma>0$ such that
\begin{equation}
\big\langle \mathbb{E}[\boldsymbol{M}_{k}(\boldsymbol{m}_{k})-\boldsymbol{M}_{k}(\dot{\boldsymbol{\omega}}_{t})],\,\boldsymbol{e}_{k}\big\rangle
\ge \frac{\varsigma}{2}\,\|\boldsymbol{e}_{k}\|_2^2.
\label{eq:appB_strong_mono}
\end{equation}
The last term in \eqref{eq:appB_add_subtract} is more challenging. In our previous works, the fixed point satisfies
$\mathbb{E}[\boldsymbol{M}_{k}(\dot{\boldsymbol{\omega}}_{t})]=\boldsymbol{0}$.
However, owing to the use of the mixed offline/online data, $\mathbb{E}[\boldsymbol{M}_{k}(\dot{\boldsymbol{\omega}}_{t})]$ is no longer exactly zero.
We decompose the surrogate mixed TD gradient as
\begin{equation}
\begin{aligned}
\boldsymbol{M}_{k}^{\mathrm{mix}}(\boldsymbol{\omega})
&\triangleq (1-\xi_{k+1})\,\boldsymbol{M}_{k}^{\mathrm{on}}(\boldsymbol{\omega})\\
&\quad+\xi_{k+1}\,\boldsymbol{M}_{k}^{\mathrm{off}}(\boldsymbol{\omega}),\\
\boldsymbol{M}_{k}(\cdot)
&\equiv \boldsymbol{M}_{k}^{\mathrm{mix}}(\cdot).
\end{aligned}
\label{eq:appB_mix_decomp}
\end{equation}
Since $\dot{\boldsymbol{\omega}}_{t}$ is the projection fixed point of the frozen online operator (for $k\in\{n_t+1,\ldots,t\}$), we have
$\mathbb{E}[\boldsymbol{M}_{k}^{\mathrm{on}}(\dot{\boldsymbol{\omega}}_{t})]=\boldsymbol{0}$. Hence,
\begin{equation}
\begin{aligned}
\mathbb{E}[\boldsymbol{M}_{k}(\dot{\boldsymbol{\omega}}_{t})]
&=(1-\xi_{k+1})\underbrace{\mathbb{E}[\boldsymbol{M}_{k}^{\mathrm{on}}(\dot{\boldsymbol{\omega}}_{t})]}_{=\boldsymbol{0}}
+\xi_{k+1}\,\mathbb{E}[\boldsymbol{M}_{k}^{\mathrm{off}}(\dot{\boldsymbol{\omega}}_{t})]\\
&=\xi_{k+1}\,\boldsymbol{b}_{\mathrm{off}}(t),
\qquad
\boldsymbol{b}_{\mathrm{off}}(t)\triangleq \mathbb{E}[\boldsymbol{M}_{k}^{\mathrm{off}}(\dot{\boldsymbol{\omega}}_{t})].
\end{aligned}
\label{eq:appB_offline_drift}
\end{equation}
By boundedness of the gradient in Appendix \ref{sec:Technical-Bounds-DNN}, there exists a constant $a_{11}>0$ such that
\begin{equation}
\|\boldsymbol{b}_{\mathrm{off}}(t)\|_2\le a_{11}m_{Q}^{1/2}.
\label{eq:appB_boff_bound}
\end{equation}
Applying Young's inequality to the drift inner product yields, for any $\varsigma>0$,
\begin{flalign}
&2\eta_{k+1}\big|\langle \mathbb{E}[\boldsymbol{M}_{k}(\dot{\boldsymbol{\omega}}_{t})],\,\boldsymbol{e}_{k}\rangle\big| &&\nonumber\\
&\le \eta_{k+1}\frac{\varsigma}{2}\,\|\boldsymbol{e}_{k}\|_2^2+\frac{2\eta_{k+1}}{\varsigma}\,\bigl\|\mathbb{E}[\boldsymbol{M}_{k}(\dot{\boldsymbol{\omega}}_{t})]\bigr\|_2^2 &&\nonumber\\
&=\eta_{k+1}\frac{\varsigma}{2}\,\|\boldsymbol{e}_{k}\|_2^2+\mathcal{O}\!\left(\eta_{k+1}\,m_Q\,\xi_{n_t}^{2}\right). &&\label{eq:appB_young_drift}
\end{flalign}

Substituting \eqref{eq:appB_add_subtract}--\eqref{eq:appB_young_drift} into \eqref{eq:appB_recursion_raw}: 
\begin{equation}
\begin{aligned}
\eta_{k+1}\|\boldsymbol{e}_{k}\|_2^2
&\le \mathcal{O}(\|\boldsymbol{e}_{k}\|_2^2-\|\boldsymbol{e}_{k+1}\|_2^2)+\mathcal{O}(m_Q\eta_{k+1}^2)\\
&\quad+\mathcal{O}(m_Q\eta_{k+1}\xi_{n_t}^{2}).
\end{aligned}
\label{eq:appB_telescope_form}
\end{equation}
Summing \eqref{eq:appB_telescope_form} for $k=n_t+1,\ldots,t$ and using monotonicity of $\{\eta_k\}$ and $\{\xi_k\}$, Substituting into \eqref{eq:appB_jensen_bias2} derives
\begin{equation}
\begin{aligned}
\sum_{t'=n_t+1}^{t} w_{t,t'}\,\bigl\|\mathbb{E}[\boldsymbol{m}_{t'}]-\dot{\boldsymbol{\omega}}_{t}\bigr\|_2^2
&\le \frac{\gamma_{n_t}}{\eta_{n_t}}\cdot
\mathcal{O}\Bigl(
\frac{1}{m_Q}
+m_Q\,t^{\kappa_6}\eta_{n_t}^{2}\\
&\qquad
+m_Q\,t^{\kappa_6}\eta_{n_t}\xi_{n_t}^{2}
\Bigr).
\end{aligned}
\label{eq:appB_bias2_recent}
\end{equation}
Finally, the bias~2 has upper bound: 
\begin{equation}
\begin{aligned}
bias~2
&\le \mathcal{O}\Bigl(
\frac{(1-\gamma_t)^{t^{\kappa_6/2}}}{\sqrt{\gamma_t}}
+\sqrt{\tfrac{\gamma_{n_t}}{\eta_{n_t}}}\\
&\qquad
+m_Q\sqrt{\gamma_{n_t}\eta_{n_t}}\,t^{\kappa_6/2}
+m_Q\sqrt{\gamma_{n_t}}\,t^{\kappa_6/2}\,\xi_{n_t}
\Bigr).
\end{aligned}
\label{eq:appB_bias2_final}
\end{equation}

\subsection{Bounds about bias~3}
The difference $\mathbb{E}[\bar{\boldsymbol{\omega}}_{t}]-\mathbb{E}[\bar{\boldsymbol{m}}_{t}]$ can also be presented as the exponentially weighted averages: 
\begin{equation}
\begin{aligned}
\bigl\|\mathbb{E}[\bar{\boldsymbol{\omega}}_{t}]-\mathbb{E}[\bar{\boldsymbol{m}}_{t}]\bigr\|_2
&\le \mathcal{O}\!\left(\frac{\gamma_{n_t}}{\gamma_t}\right)\cdot e_{n_t}^{\mathrm{b}},\\
e_{n_t}^{\mathrm{b}}
&\triangleq \max_{k\in\{n_t+1,\ldots,t\}}\bigl\|\mathbb{E}[\boldsymbol{\omega}_{k}]-\mathbb{E}[\boldsymbol{m}_{k}]\bigr\|_2.
\end{aligned}
\label{eq:appB_bias3_reduce}
\end{equation}
Using the non-expansiveness of projection operation and the recursions formulation in \eqref{eq:TD-learning} and \eqref{eq:appB_aux_update}, it can be expanded as: 
\begin{equation}
\begin{aligned}
\bigl\|\mathbb{E}[\boldsymbol{\omega}_{k}]-\mathbb{E}[\boldsymbol{m}_{k}]\bigr\|_2
&\le \sum_{j=n_t+1}^{k}\eta_j\cdot\bigl\|\mathbb{E}[\boldsymbol{\Delta}(\boldsymbol{\omega}_{j-1})]
-\mathbb{E}[\boldsymbol{M}^{\mathrm{mix}}_{j-1}(\boldsymbol{m}_{j-1})]\bigr\|_2.
\end{aligned}
\label{eq:appB_bias3_sumgrad}
\end{equation}

Define the TD-gradient integrands
\begin{equation}
\begin{aligned}
H_f(\boldsymbol{\omega};\varepsilon)
&\triangleq \delta(\boldsymbol{\omega};\varepsilon)\,
\nabla_{\boldsymbol{\omega}}f(\boldsymbol{\omega};\boldsymbol{s},\boldsymbol{a}),\\
H_{\hat f}(\boldsymbol{m};\varepsilon)
&\triangleq \hat\delta(\boldsymbol{m};\varepsilon)\,
\nabla_{\boldsymbol{\omega}}\hat f(\boldsymbol{m};\boldsymbol{s},\boldsymbol{a}).
\end{aligned}
\end{equation}


where $\varepsilon=(\boldsymbol{s},\boldsymbol{a},\boldsymbol{s}',\boldsymbol{a}')$, and $\boldsymbol{a}'\sim\pi_{\boldsymbol{\theta}_t}(\cdot\mid\boldsymbol{s}')$.
Let $\mu_j^{\mathrm{on}}$ denote distribution of online data generated under $\pi_{\boldsymbol{\theta}_j}$, while $\tilde\mu_t^{\mathrm{on}}$ denote online data's distribution under the frozen policy $\pi_{\boldsymbol{\theta}_t}$, and let $\nu^{\mathrm{off}}$ denote the offline data distribution. Then \eqref{eq:appB_bias3_sumgrad} can be decomposed as: 
\begin{equation}
\begin{aligned}
&\bigl\|\mathbb{E}[\boldsymbol{\Delta}_{j-1}]-\mathbb{E}[\boldsymbol{M}^{\mathrm{mix}}_{j-1}(\boldsymbol{m}_{j-1})]\bigr\|_2
\\
&\le (1-\xi_j)\bigl\|\mathbb{E}_{\mu_j^{\mathrm{on}}}\!\big[H_f(\boldsymbol{\omega}_{j-1};\varepsilon)\big]
-\mathbb{E}_{\tilde\mu_t^{\mathrm{on}}}\!\big[H_f(\boldsymbol{\omega}_{j-1};\varepsilon)\big]\bigr\|_2
\\
&\quad+\bigl\|\mathbb{E}_{\tilde\mu_t^{\mathrm{on}}}\!\big[H_f(\boldsymbol{\omega}_{j-1};\varepsilon)-H_{\hat f}(\boldsymbol{m}_{j-1};\varepsilon)\big]\bigr\|_2
\\
&\quad+\xi_j\bigl\|\mathbb{E}_{\nu^{\mathrm{off}}}\!\big[H_f(\boldsymbol{\omega}_{j-1};\varepsilon)-H_{\hat f}(\boldsymbol{m}_{j-1};\varepsilon)\big]\bigr\|_2.
\label{eq:appB_grad_decomp}
\end{aligned}
\end{equation}

We further obtain
\begin{equation}
\begin{aligned}
&\bigl\|\mathbb{E}[\boldsymbol{\Delta}_{j-1}]
-\mathbb{E}[\boldsymbol{M}^{\mathrm{mix}}_{j-1}(\boldsymbol{m}_{j-1})]\bigr\|_2\\
&\le m_Q^{1/2}\cdot\mathcal{O}\!\Bigl(
\|\mu_j^{\mathrm{on}}-\tilde\mu_t^{\mathrm{on}}\|_{\mathrm{TV}}
+\|\boldsymbol{\theta}_j-\boldsymbol{\theta}_t\|_2
\Bigr)\\
&\quad+\mathcal{O}(\epsilon_{m_Q}).
\end{aligned}
\label{eq:appB_grad_tv}
\end{equation}

Following the ergodicity assumption, 
\begin{equation}
\|\mu_j^{\mathrm{on}}-\tilde\mu_t^{\mathrm{on}}\|_{\mathrm{TV}}
\le \lambda\chi^{\tau_j}
+\mathcal{O}\!\left(\sum_{k=j-\tau_j+1}^{t}\beta_k\right),
\label{eq:appB_tv_bound}
\end{equation}
where $\tau_j$ is the mixing-time parameter and $\chi^{\tau_j}=\mathcal{O}(1/t)$. Moreover, by Lipschitz continuity of $\pi_{\boldsymbol{\theta}}$,
\begin{equation}
\|\boldsymbol{\theta}_j-\boldsymbol{\theta}_t\|_2
\le \sum_{k=j+1}^{t}\beta_k
\le \mathcal{O}\!\bigl(\beta_{n_t}t^{\kappa_6}\bigr),
\label{eq:appB_theta_drift}
\end{equation}
where the last step uses $j\ge n_t+1=t-\Theta(t^{\kappa_6})$. Substituting \eqref{eq:appB_tv_bound}--\eqref{eq:appB_theta_drift} into \eqref{eq:appB_grad_tv}, and then into \eqref{eq:appB_bias3_sumgrad}, yields
\begin{equation}
\begin{aligned}
e_{n_t}^{\mathrm{b}}
&\le m_Q^{1/2}\eta_{n_t}\cdot\mathcal{O}\!\Bigl(
t^{\kappa_6-1}+\beta_{n_t}t^{2\kappa_6}
\Bigr)
+\mathcal{O}(\eta_{n_t}\epsilon_{m_Q}t^{\kappa_6}).
\end{aligned}
\label{eq:appB_ebound}
\end{equation}

Substituting \eqref{eq:appB_ebound} into \eqref{eq:appB_bias3_reduce}, the bias 3 can be bounded as: 
\begin{equation}
\begin{aligned}
Bias~3
&\le \mathcal{O}\Bigl(
m_Q\eta_{n_t}t^{\kappa_6-1}
+m_Q\eta_{n_t}\beta_{n_t}t^{2\kappa_6}
\Bigr)\\
&\quad+\mathcal{O}(\epsilon_{m_Q}).
\end{aligned}
\label{eq:appB_bias3_final}
\end{equation}

Substituting \eqref{eq:appB_bias2_final} and \eqref{eq:appB_bias3_final} into \eqref{eq:appB_core_goal}, the upper bound of the critic error can be presented as: 
\begin{equation}
\begin{aligned}
\epsilon_{t}^{\mathrm{cri}}
&\le \mathcal{O}\Bigl(
\epsilon_{m_Q}
+\frac{(1-\gamma_{t})^{t^{\kappa_{6}/2}}}{\sqrt{\gamma_{t}}}
+\sqrt{\tfrac{\gamma_{n_{t}}}{\eta_{n_{t}}}}\\
&\qquad
+m_Q\sqrt{\gamma_{n_{t}}\eta_{n_{t}}}\,t^{\kappa_{6}/2}
+m_Q\eta_{n_{t}}\,t^{\kappa_{6}-1}
+m_Q\eta_{n_{t}}\beta_{n_{t}}\,t^{2\kappa_{6}}\\
&\qquad
+m_Q\sqrt{\gamma_{n_t}}\,t^{\kappa_6/2}\,\xi_{n_t}
\Bigr),
\end{aligned}
\end{equation}
which is presented in Lemma~\ref{lemma-convergence-rate-critic}. This completes the proof.  
\section{Proof of Lemma \ref{lemma-asymptotic-consistency}}
\label{sec:proof-of-asymptotic-consistency}
In this section, we restate the technical lemma in [xxx]. 
\begin{lemma}[A projected stochastic approximation tracking lemma]
\label{lemma:appC_SA_tracking}
Consider a filtration $\{\mathcal{F}_t\}_{t\ge 0}$ and the recursion
\begin{equation}
\boldsymbol{w}^{t+1}=\Pi_{\mathcal{W}}\!\Big(\boldsymbol{w}^{t}+\alpha_t\big(\boldsymbol{\varrho}^{t+1}-\boldsymbol{w}^{t}\big)\Big),
\label{eq:appC_SA_recursion}
\end{equation}
where $\Pi_{\mathcal{W}}$ is the Euclidean projection onto a nonempty, compact, and convex set $\mathcal{W}$.
Assume:
\begin{enumerate}
\item $\boldsymbol{w}^{t}$ is $\mathcal{F}_t$-measurable and $\boldsymbol{\varrho}^{t+1}$ is $\mathcal{F}_{t+1}$-measurable;
\item there exists a constant $C<\infty$ such that $\mathbb{E}[\|\boldsymbol{\varrho}^{t+1}\|_2^2\mid\mathcal{F}_t]\le C$ a.s. for all $t$;
\item there exist an $\mathcal{F}_t$-measurable process $\{\boldsymbol{z}^t\}$ and an $\mathcal{F}_t$-measurable bias term $\{\boldsymbol{o}^t\}$ such that
\begin{equation}
\mathbb{E}[\boldsymbol{\varrho}^{t+1}\mid\mathcal{F}_t]=\boldsymbol{z}^t+\boldsymbol{o}^t,\qquad \sum_{t=0}^{\infty}\alpha_t\,\|\boldsymbol{o}^t\|_2<\infty\ \ \text{a.s.};
\label{eq:appC_cond_mean}
\end{equation}
\item the step size satisfies $\alpha_t>0$, $\sum_t\alpha_t=\infty$, and $\sum_t\alpha_t^2<\infty$;
\item the target drifts slowly: $\|\boldsymbol{z}^{t+1}-\boldsymbol{z}^t\|_2/\alpha_t\to 0$ a.s.
\end{enumerate}
Then $\|\boldsymbol{w}^{t}-\boldsymbol{z}^t\|_2\to 0$ almost surely.
\end{lemma}

\subsection{Asymptotic consistency of $\hat{J}_i^t$}
Rewrite the recursion \eqref{eq:J-average} as
\begin{equation}
\hat{J}_i^{t}=\hat{J}_i^{t-1}+\alpha_t\big(\tilde{J}_i^{t}-\hat{J}_i^{t-1}\big).
\label{eq:appC_J_recursion}
\end{equation}
$\tilde{J}_i^{t}$ can be presented as the combination of offline and online part: 
\begin{equation}
\tilde{J}_i^{t}=(1-\xi_t)\,\hat{\mathbb{E}}_{\mathcal{D}_0^{t}}\big[C_i^{\text{'}}(\boldsymbol{s},\boldsymbol{a})\big]
+\xi_t\,\hat{\mathbb{E}}_{\mathcal{D}_{\mathrm{off}}^{t}}\big[C_i^{\text{'}}(\boldsymbol{s},\boldsymbol{a})\big].
\end{equation}
Taking conditional expectation given $\mathcal{F}_{t-1}$ yields
\begin{equation}
\mathbb{E}[\tilde{J}_i^{t}\mid\mathcal{F}_{t-1}]=J_i(\boldsymbol{\theta}_{t-1})+o_{J}^{t-1},
\label{eq:appC_J_cond_mean}
\end{equation}
where
\begin{equation}
J_i(\boldsymbol{\theta})\triangleq \mathbb{E}_{\sigma_{\pi_{\boldsymbol{\theta}}}}\big[C_i^{\text{'}}(\boldsymbol{s},\boldsymbol{a})\big]
\label{eq:appC_J_def}
\end{equation}
is the stationary average cost under $\pi_{\boldsymbol{\theta}}$ and $o_{J}^{t-1}$ denotes the bias term to be bounded. 

Let $\mu_k$ denote the conditional distribution of $(\boldsymbol{s}_k,\boldsymbol{a}_k)$ given $\mathcal{F}_{t-1}$ for $k\in\{t-T_t+1,\ldots,t\}$. 
The online bias can be written as
\begin{equation}
\begin{aligned}
 |b_{J,\mathrm{on}}^{t-1}|
 &\triangleq
 |\mathbb{E}\Big[\hat{\mathbb{E}}_{\mathcal{D}_0^{t}}[C_i^{\text{'}}]\,\Big|\,\mathcal{F}_{t-1}\Big]
 - J_i(\boldsymbol{\theta}_{t-1})| \\
 &= \bigg|
 \frac{1}{T_t}\sum_{k=t-T_t+1}^{t}
 \Big(\mathbb{E}_{\mu_k}[C_i^{\text{'}}]-\mathbb{E}_{\sigma_{\pi_{\boldsymbol{\theta}_{t-1}}}}[C_i^{\text{'}}]\Big)
 \bigg| \\
 &\le
 \frac{1}{T_t}\sum_{k=t-T_t+1}^{t}
 \Big|
 \mathbb{E}_{\mu_k}[C_i^{\text{'}}]-\mathbb{E}_{\sigma_{\pi_{\boldsymbol{\theta}_{t-1}}}}[C_i^{\text{'}}]
 \Big| \\
 &\le
 \frac{2C_{\max}}{T_t}\sum_{k=t-T_t+1}^{t}
 \big\|\mu_k-\sigma_{\pi_{\boldsymbol{\theta}_{t-1}}}\big\|_{\mathrm{TV}}.
\end{aligned}
\label{eq:appC_bJon_tv}
\end{equation}

For each $k$ in the window, apply the triangle inequality
\begin{equation}
\big\|\mu_k-\sigma_{\pi_{\boldsymbol{\theta}_{t-1}}}\big\|_{\mathrm{TV}}
\le
\underbrace{\big\|\mu_k-\sigma_{\pi_{\boldsymbol{\theta}_{k}}}\big\|_{\mathrm{TV}}}_{\text{mixing term}}
+
\underbrace{\big\|\sigma_{\pi_{\boldsymbol{\theta}_{k}}}-\sigma_{\pi_{\boldsymbol{\theta}_{t-1}}}\big\|_{\mathrm{TV}}}_{\text{policy-drift term}}.
\label{eq:appC_tv_split}
\end{equation}
The mixing term is controlled by Assumption~\ref{assumption:problem_structure}, for some constants $\lambda>0$ and $\varrho\in(0,1)$, $\|\mu_k-\sigma_{\pi_{\boldsymbol{\theta}_k}}\|_{\mathrm{TV}}\le \lambda\varrho^{\tau_t}$ with a mixing horizon $\tau_t=\Theta(\log t)$. Since $T_t=\mathcal{O}(\log t)$, we can take $\tau_t\asymp T_t$, so that $\varrho^{\tau_t}=\mathcal{O}(t^{-2})$. For the policy-drift term, Assumption~\ref{assumption:problem_structure} implies Lipschitz continuity of the stationary distribution in total variation:
there exists $L_{\sigma}>0$ such that
$\|\sigma_{\pi_{\boldsymbol{\theta}}}-\sigma_{\pi_{\boldsymbol{\theta}'}}\|_{\mathrm{TV}}\le L_{\sigma}\|\boldsymbol{\theta}-\boldsymbol{\theta}'\|_2$. Moreover, using the compactness of $\boldsymbol{\Theta}$, there exists $D_{\Theta}<\infty$ such that
\begin{equation}
\|\boldsymbol{\theta}_{\ell+1}-\boldsymbol{\theta}_{\ell}\|_2\le D_{\Theta}\,\beta_{\ell},\qquad \forall \ell.
\label{eq:appC_theta_step}
\end{equation}
Hence for $k\in\{t-T_t+1,\ldots,t\}$,
\begin{equation}
\|\boldsymbol{\theta}_{k}-\boldsymbol{\theta}_{t-1}\|_2
\le \sum_{\ell=k}^{t-2}\|\boldsymbol{\theta}_{\ell+1}-\boldsymbol{\theta}_{\ell}\|_2
\le D_{\Theta}\sum_{\ell=t-T_t}^{t-1}\beta_{\ell}
\le D_{\Theta}T_t\,\beta_{t-T_t+1}.
\label{eq:appC_theta_drift_window}
\end{equation}
Then the online bias can be bounded as: 
\begin{equation}
|b_{J,\mathrm{on}}^{t-1}|
\le c_{J,1}\Big(\varrho^{\tau_t}+T_t\,\beta_{t-T_t+1}\Big)
\le c_{J,1}\Big(t^{-2}+T_t\,\beta_{t-T_t+1}\Big),
\label{eq:appC_bJon_bound}
\end{equation}
for some constant $c_{J,1}>0$.

For the offline part, with the boundedness of $C_i^{\text{'}}$,
\begin{equation}
\bigg|\mathbb{E}\Big[\hat{\mathbb{E}}_{\mathcal{D}_{\mathrm{off}}^{t}}[C_i^{\text{'}}]\,\Big|\,\mathcal{F}_{t-1}\Big]-J_i(\boldsymbol{\theta}_{t-1})\bigg|
\le 2C_{\max}
\triangleq c_{J,2}.
\label{eq:appC_bJoff_bound}
\end{equation}
Therefore, combining \eqref{eq:appC_bJon_bound} and \eqref{eq:appC_bJoff_bound}, the total bias in \eqref{eq:appC_J_cond_mean} can be bounded as: 
\begin{equation}
|o_J^{t-1}|
\le (1-\xi_t)|b_{J,\mathrm{on}}^{t-1}|+\xi_t\,c_{J,2}
\le c_{J,1}\Big(t^{-2}+T_t\,\beta_{t-T_t+1}\Big)+c_{J,2}\,\xi_t.
\label{eq:appC_oJ_bound}
\end{equation}

The boundedness condition in Lemma~\ref{lemma:appC_SA_tracking}.2 holds because $|\tilde J_i^t|\le C_{\max}$.
By Assumption~\ref{assumption:step_size}.1 and the summability requirements in Assumption~\ref{assumption:step_size}.2 and Assumption~\ref{assumption:step_size}.5, together with $T_t=\mathcal{O}(\log t)$, the bias bound \eqref{eq:appC_oJ_bound} implies $\sum_t \alpha_t |o_J^{t-1}|<\infty$. Also, $J_i(\boldsymbol{\theta})$ is Lipschitz over the compact set $\boldsymbol{\Theta}$, $|J_i(\boldsymbol{\theta}_{t})-J_i(\boldsymbol{\theta}_{t-1})|\le L_J\|\boldsymbol{\theta}_{t}-\boldsymbol{\theta}_{t-1}\|_2\le L_J D_{\Theta}\beta_{t-1}$. 
Thus $|J_i(\boldsymbol{\theta}_{t})-J_i(\boldsymbol{\theta}_{t-1})|/\alpha_{t-1}\to 0$ by Assumption~\ref{assumption:step_size}.2 (i.e., $\beta_t/\alpha_t\to 0$), which verifies Lemma~\ref{lemma:appC_SA_tracking}.5.

Applying Lemma~\ref{lemma:appC_SA_tracking} to \eqref{eq:appC_J_recursion} with
$\boldsymbol{w}^{t}=\hat J_i^{t}$,
$\boldsymbol{\varrho}^{t}=\tilde J_i^{t}$, and
$\boldsymbol{z}^{t-1}=J_i(\boldsymbol{\theta}_{t-1})$,
we obtain
\begin{equation}
|\hat J_i^{t}-J_i(\boldsymbol{\theta}_{t-1})|\to 0,\qquad \text{a.s.}
\end{equation}
Finally, since $|J_i(\boldsymbol{\theta}_{t})-J_i(\boldsymbol{\theta}_{t-1})|\to 0$, the above implies
\begin{equation}
|\hat J_i^{t}-J_i(\boldsymbol{\theta}_{t})|\to 0,\qquad \text{a.s.}
\label{eq:appC_J_consistency}
\end{equation}

\subsection{Asymptotic consistency of $\hat{\boldsymbol{g}}_i^t$ up to $\epsilon_{m_Q}$}

Define the gradient of the mixed policy  score function as $\boldsymbol{\phi}_{\boldsymbol{\theta}}(\boldsymbol{s},\boldsymbol{a})\triangleq\nabla_{\boldsymbol{\theta}}\log\pi_{\boldsymbol{\theta}}(\boldsymbol{a}\mid\boldsymbol{s})$. 
Under Assumption~\ref{assumption:problem_structure} and the mixture-ratio terms are uniformly bounded over $\boldsymbol{\Theta}$, there exists a constant $B_{\pi}<\infty$ such that
\begin{equation}
\sup_{\boldsymbol{\theta}\in\boldsymbol{\Theta}}\ \sup_{(\boldsymbol{s},\boldsymbol{a})\in\mathcal{S}\times\mathcal{A}}
\big\|\boldsymbol{\phi}_{\boldsymbol{\theta}}(\boldsymbol{s},\boldsymbol{a})\big\|_2\le B_{\pi}.
\label{eq:appC_score_bound}
\end{equation}

Then, define an auxiliary gradient target: 
\begin{equation}
\nabla_{\boldsymbol{\theta}}\hat{J}_i(\boldsymbol{\theta}_t)
\triangleq
\mathbb{E}_{\sigma_{\pi_{\boldsymbol{\theta}_t}}}\Big[\hat{Q}_i^{\pi_{\boldsymbol{\theta}_t}}(\boldsymbol{s},\boldsymbol{a})\,\boldsymbol{\phi}_{\boldsymbol{\theta}_t}(\boldsymbol{s},\boldsymbol{a})\Big]. 
\label{eq:appC_aux_grad_def}
\end{equation}
$\hat{\boldsymbol{g}}_i^t$ tracks the moving target $\nabla_{\boldsymbol{\theta}}\hat{J}_i(\boldsymbol{\theta}_t)$ as
\begin{equation}
\hat{\boldsymbol{g}}_i^{t}=\hat{\boldsymbol{g}}_i^{t-1}+\alpha_t\big(\tilde{\boldsymbol{g}}_i^{t}-\hat{\boldsymbol{g}}_i^{t-1}\big).
\label{eq:appC_g_recursion}
\end{equation}

From conditional expectation of \eqref{eq:gtlt} given $\mathcal{F}_{t-1}$, we can derive: 
\begin{equation}
\mathbb{E}[\tilde{\boldsymbol{g}}_i^{t}\mid\mathcal{F}_{t-1}]
=\nabla_{\boldsymbol{\theta}}\hat{J}_i(\boldsymbol{\theta}_{t-1})+\boldsymbol{o}_g^{t-1},
\label{eq:appC_g_cond_mean}
\end{equation}
where $\boldsymbol{o}_g^{t-1}$ is the bias to be bounded. It can also be decomposed into offline and online part. Adding and subtracting the ideal target $\nabla_{\boldsymbol{\theta}}\hat{J}_i(\boldsymbol{\theta}_{t-1})$, we can obtain: 
\begin{equation}
\begin{aligned}
\boldsymbol{o}_g^{t-1}
=
&(1-\xi_t)\Big(
\mathbb{E}[\tilde{\boldsymbol{g}}_{i,\mathrm{on}}^t\mid\mathcal{F}_{t-1}]
-\nabla_{\boldsymbol{\theta}}\hat{J}_i(\boldsymbol{\theta}_{t-1})
\Big) \\
&+\xi_t\Big(
\mathbb{E}[\tilde{\boldsymbol{g}}_{i,\mathrm{off}}^t\mid\mathcal{F}_{t-1}]
-\nabla_{\boldsymbol{\theta}}\hat{J}_i(\boldsymbol{\theta}_{t-1})
\Big).
\end{aligned}
\end{equation}
The online part can be decomposed into two errors: critic error and distribution difference. It can be bounded by the same method of our previous work SLDAC. The offline part of bias can be bounded using the similar trick in \eqref{eq:appC_bJoff_bound} under the boundedness of Q function and score function. Then, the bias has upper bound as: 
\begin{equation}
\|\boldsymbol{o}_{g}^{t-1}\|_2
\le c_{g,1}\Big(t^{-2}+T_t\,\beta_{t-T_t+1}\Big)
+c_{g,2}\,\xi_t
+ c_{g,3}\,\bar{\epsilon}_{\mathrm{cri}}(t),
\label{eq:appC_og_bound}
\end{equation}
where $\bar{\epsilon}_{\mathrm{cri}}(t)\triangleq\max_{0\le j\le I}\epsilon_{t,j}^{\mathrm{cri}}$ is the aggregated critic tracking error and $c_{g,1},c_{g,2},c_{g,3}>0$ are constants independent of $t$ and $m_Q$.
Combining with Assumption~\ref{assumption:step_size}, $\sum_t \alpha_t\|\boldsymbol{o}_g^{t-1}\|_2<\infty$ is satisfied. Moreover, since $\nabla_{\boldsymbol{\theta}}\hat{J}_i(\boldsymbol{\theta})$ is Lipschitz over the compact $\boldsymbol{\Theta}$ and $\|\boldsymbol{\theta}_{t}-\boldsymbol{\theta}_{t-1}\|_2=\mathcal{O}(\beta_{t-1})$ by \eqref{eq:appC_theta_step}, we can obtain that $\|\nabla_{\boldsymbol{\theta}}\hat{J}_i(\boldsymbol{\theta}_{t})-\nabla_{\boldsymbol{\theta}}\hat{J}_i(\boldsymbol{\theta}_{t-1})\|_2/\alpha_{t-1}\to 0$ when $\beta_t/\alpha_t\to 0$. In this case, all conditions of Lemma~\ref{lemma:appC_SA_tracking} are satisfied so that: 
\begin{equation}
\big\|\hat{\boldsymbol{g}}_i^{t}-\nabla_{\boldsymbol{\theta}}\hat{J}_i(\boldsymbol{\theta}_{t})\big\|_2\to 0,
\qquad \text{a.s.}
\label{eq:appC_g_track_aux}
\end{equation}

The bias between auxiliary policy gradient estimate and true policy gradient can be presented as: 
\begin{align}
\big\|\nabla_{\boldsymbol{\theta}}\hat{J}_i(\boldsymbol{\theta}_t)-\nabla_{\boldsymbol{\theta}}J_i(\boldsymbol{\theta}_t)\big\|_2
&=\Big\|\mathbb{E}_{\sigma_{\pi_{\boldsymbol{\theta}_t}}}\big[\big(\hat{Q}_i^{\pi_{\boldsymbol{\theta}_t}}-Q_i^{\pi_{\boldsymbol{\theta}_t}}\big)\,\boldsymbol{\phi}_{\boldsymbol{\theta}_t}\big]\Big\|_2 \nonumber\\
&\le B_{\pi}\,\mathbb{E}_{\sigma_{\pi_{\boldsymbol{\theta}_t}}}\Big[\big|\hat{Q}_i^{\pi_{\boldsymbol{\theta}_t}}(\boldsymbol{s},\boldsymbol{a})-Q_i^{\pi_{\boldsymbol{\theta}_t}}(\boldsymbol{s},\boldsymbol{a})\big|\Big].
\label{eq:appC_grad_diff_reduce}
\end{align}
Under the local-linearization analysis of the critic in Appendix \ref{sec:proof-of-critic} and Assumption~\ref{assumption:target_Q}, $\mathbb{E}_{\sigma_{\pi_{\boldsymbol{\theta}_t}}}\Big[\big|\hat{Q}_i^{\pi_{\boldsymbol{\theta}_t}}-Q_i^{\pi_{\boldsymbol{\theta}_t}}\big|\Big]$ admits the uniform bound: 
\begin{equation}
\mathbb{E}_{\sigma_{\pi_{\boldsymbol{\theta}_t}}}\Big[\big|\hat{Q}_i^{\pi_{\boldsymbol{\theta}_t}}-Q_i^{\pi_{\boldsymbol{\theta}_t}}\big|\Big]
\le c_Q\Big(|\hat{J}_i^t-J_i(\boldsymbol{\theta}_t)|+\epsilon_{m_Q}\Big),
\label{eq:appC_Q_mismatch_bound}
\end{equation}
for some constant $c_Q>0$ independent of $t$ and $m_Q$.
Combining \eqref{eq:appC_grad_diff_reduce}--\eqref{eq:appC_Q_mismatch_bound} and using \eqref{eq:appC_J_consistency} yields
\begin{equation}
\limsup_{t\to\infty}\big\|\nabla_{\boldsymbol{\theta}}\hat{J}_i(\boldsymbol{\theta}_t)-\nabla_{\boldsymbol{\theta}}J_i(\boldsymbol{\theta}_t)\big\|_2
\le c_{\pi}\,\epsilon_{m_Q},
\label{eq:appC_aux_to_true_grad}
\end{equation}
where $c_{\pi}\triangleq B_{\pi}c_Q$.

Using the triangle inequality,
\begin{equation}
\begin{aligned}
\big\|\hat{\boldsymbol{g}}_i^{t}-\nabla_{\boldsymbol{\theta}}J_i(\boldsymbol{\theta}_t)\big\|_2
&\le \big\|\hat{\boldsymbol{g}}_i^{t}-\nabla_{\boldsymbol{\theta}}\hat{J}_i(\boldsymbol{\theta}_t)\big\|_2 \\
&\quad + \big\|\nabla_{\boldsymbol{\theta}}\hat{J}_i(\boldsymbol{\theta}_t)-\nabla_{\boldsymbol{\theta}}J_i(\boldsymbol{\theta}_t)\big\|_2.
\end{aligned}
\label{eq:asm2}
\end{equation}
When $t\to\infty$, the first term vanishes and the second term is bounded by $c_{\pi}\epsilon_{m_Q}$ by \eqref{eq:appC_aux_to_true_grad}.
Combining \eqref{eq:asm2} with \eqref{eq:appC_J_consistency}, we complete the proof of Lemma~\ref{lemma-asymptotic-consistency}. 

\section{Proof of Lemma \ref{lemma-convergence-rate-surrogate}}
\label{sec:proof-of-surrogate}
Since the analysis for $\epsilon_{J}(t)$ and $\epsilon_{g}(t)$ are similar, we only provide the more complicated derivation of the latter due to space limit. 

Firstly, we define the accumulated gradient-tracking error
\begin{equation}
\begin{aligned}
\boldsymbol{y}_k
&\triangleq \hat{\boldsymbol{g}}_i^k
-\nabla_{\boldsymbol{\theta}}J_i(\boldsymbol{\theta}_k), \\
\epsilon_{g,i}(t)
&\triangleq \frac{1}{t+1}\sum_{k=0}^{t}
\mathbb{E}\big[\|\boldsymbol{y}_k\|_2^2\big].
\end{aligned}
\label{eq:appD_epsgi_def}
\end{equation}
Let
\begin{equation}
 \boldsymbol{\Delta}_k\triangleq \nabla_{\boldsymbol{\theta}}J_i(\boldsymbol{\theta}_k)-\nabla_{\boldsymbol{\theta}}J_i(\boldsymbol{\theta}_{k+1})
\label{eq:appD_Deltag_def}
\end{equation}
and define the cross term
\begin{equation}
 \Xi_{k+1}\triangleq
 \big\langle
 \boldsymbol{y}_k,
 \tilde{\boldsymbol{g}}_i^{k+1}-\nabla_{\boldsymbol{\theta}}J_i(\boldsymbol{\theta}_k)
 \big\rangle.
\label{eq:appD_Xi_def}
\end{equation}
According to \eqref{eq:g-average}, we have the following recursion for $\boldsymbol{y}_k$: 
\begin{align}
 \boldsymbol{y}_{k+1}
 &=\hat{\boldsymbol{g}}_i^{k+1}-\nabla_{\boldsymbol{\theta}}J_i(\boldsymbol{\theta}_{k+1}) \notag\\
 &=\boldsymbol{y}_k+\boldsymbol{\Delta}_k+\alpha_{k+1}\big(\tilde{\boldsymbol{g}}_i^{k+1}-\hat{\boldsymbol{g}}_i^k\big).
\label{eq:appD_y_recursion}
\end{align}
Using $2\langle a,b\rangle\le \|a\|_2^2+\|b\|_2^2$ and noting that
\begin{equation}
 \big\langle \boldsymbol{y}_k,\tilde{\boldsymbol{g}}_i^{k+1}-\hat{\boldsymbol{g}}_i^k\big\rangle
 =\Xi_{k+1}-\|\boldsymbol{y}_k\|_2^2,
\end{equation}
we obtain
\begin{equation}
\begin{aligned}
 \|\boldsymbol{y}_{k+1}\|_2^2
 &\le (1-2\alpha_{k+1})\|\boldsymbol{y}_k\|_2^2
 +2\alpha_{k+1}\Xi_{k+1} \\
 &\quad +2\langle \boldsymbol{y}_k,\boldsymbol{\Delta}_k\rangle
 +2\|\boldsymbol{\Delta}_k\|_2^2 \\
 &\quad +2\alpha_{k+1}^2
 \big\|\tilde{\boldsymbol{g}}_i^{k+1}-\hat{\boldsymbol{g}}_i^k\big\|_2^2.
\end{aligned}
\label{eq:appD_y_square}
\end{equation}
It can be decomposed into five parts: 
\begin{align}
 \epsilon_{g,i}(t)
 &\le
 \underbrace{\frac{1}{t+1}\sum_{k=0}^{t}\frac{\mathbb{E}[\|\boldsymbol{y}_k\|_2^2-\|\boldsymbol{y}_{k+1}\|_2^2]}{2\alpha_{k+1}}}_{M_1(t)}
 +\underbrace{\frac{1}{t+1}\sum_{k=0}^{t}\frac{1}{\alpha_{k+1}}\mathbb{E}[\|\boldsymbol{\Delta}_k\|_2^2]}_{M_2(t)} \notag\\
 &\quad +\underbrace{\frac{1}{t+1}\sum_{k=0}^{t}\alpha_{k+1}\mathbb{E}\Big[\big\|\tilde{\boldsymbol{g}}_i^{k+1}-\hat{\boldsymbol{g}}_i^k\big\|_2^2\Big]}_{M_3(t)}
 +\underbrace{\frac{1}{t+1}\sum_{k=0}^{t}\mathbb{E}[\Xi_{k+1}]}_{M_4(t)} \notag\\
 &\quad +\underbrace{\frac{1}{t+1}\sum_{k=0}^{t}\frac{1}{\alpha_{k+1}}\mathbb{E}[\langle \boldsymbol{y}_k,\boldsymbol{\Delta}_k\rangle]}_{M_5(t)}.
\label{eq:appD_gradient_decomp}
\end{align}
In the following, we bound five terms separately. 

For $M_1(t)$, 
\begin{equation}
\begin{aligned}
\sum_{k=0}^{t}
\frac{\mathbb{E}[\|\boldsymbol{y}_k\|_2^2-\|\boldsymbol{y}_{k+1}\|_2^2]}
{2\alpha_{k+1}}
={}& \frac{\mathbb{E}[\|\boldsymbol{y}_0\|_2^2]}{2\alpha_1}
-\frac{\mathbb{E}[\|\boldsymbol{y}_{t+1}\|_2^2]}{2\alpha_{t+1}} \\
&\quad +\sum_{k=1}^{t}\left(
\frac{1}{2\alpha_{k+1}}-\frac{1}{2\alpha_k}\right)
\mathbb{E}[\|\boldsymbol{y}_k\|_2^2].
\end{aligned}
\end{equation}
since $\alpha_k$ is non-increasing and $\mathbb{E}[\|\boldsymbol{y}_k\|_2^2]$ is uniformly bounded, we have
\begin{equation}
 M_1(t)=\mathcal{O}\!\Big(\frac{1}{(t+1)\alpha_{t+1}}\Big),
\label{eq:appD_M1_bound}
\end{equation}

For $M_2(t)$, $\nabla_{\boldsymbol{\theta}}J_i(\cdot)$ is Lipschitz on the compact set $\boldsymbol{\Theta}$, there exists $L_i>0$ such that

\begin{equation}
\|\boldsymbol{\Delta}_k\|_2
=\big\|\nabla_{\boldsymbol{\theta}}J_i(\boldsymbol{\theta}_k)-\nabla_{\boldsymbol{\theta}}J_i(\boldsymbol{\theta}_{k+1})\big\|_2
\le L_i\|\boldsymbol{\theta}_{k+1}-\boldsymbol{\theta}_k\|_2.
\end{equation}
The actor update \eqref{eq:theta update} gives
$
\|\boldsymbol{\theta}_{k+1}-\boldsymbol{\theta}_k\|_2
=\beta_k\|\bar{\boldsymbol{\theta}}_k-\boldsymbol{\theta}_k\|_2
\le c\beta_k,
$
with the compact set $\boldsymbol{\Theta}$. Hence
\begin{equation}
\mathbb{E}[\|\boldsymbol{\Delta}_k\|_2^2]\le c\beta_k^2,
\end{equation}
which implies
\begin{equation}
M_2(t)=\mathcal{O}\!\Bigg(\frac{1}{t+1}\sum_{k=0}^{t}\frac{\beta_k^2}{\alpha_{k+1}}\Bigg).
\label{eq:appD_M2_bound}
\end{equation}

For $M_3(t)$, $\tilde{\boldsymbol{g}}_i^{k+1}$ is formed from the bounded score function and the bounded critic output. There exists a constant $G_2<\infty$ such that
\begin{equation}
\sup_k\mathbb{E}\big[\|\tilde{\boldsymbol{g}}_i^{k+1}\|_2^2\big]\le G_2.
\end{equation}
Using $\|a-b\|_2^2\le 2\|a\|_2^2+2\|b\|_2^2$, we obtain
\begin{equation}
\mathbb{E}\Big[\big\|\tilde{\boldsymbol{g}}_i^{k+1}-\hat{\boldsymbol{g}}_i^k\big\|_2^2\Big]\le c,
\end{equation}
and therefore
\begin{equation}
M_3(t)=\mathcal{O}\!\Bigg(\frac{1}{t+1}\sum_{k=0}^{t}\alpha_{k+1}\Bigg).
\label{eq:appD_M3_bound}
\end{equation}

For $M_4(t)$, according to Appendix~\ref{sec:proof-of-asymptotic-consistency}, the conditional mean of the instantaneous gradient estimate satisfies
\begin{equation}
\mathbb{E}[\tilde{\boldsymbol{g}}_i^{k+1}\mid \mathcal{F}_k]
=\nabla_{\boldsymbol{\theta}}\hat{J}_i(\boldsymbol{\theta}_k)+\boldsymbol{o}_{g,i}^k,
\label{eq:appD_conditional_mean}
\end{equation}
where
\begin{equation}
\|\boldsymbol{o}_{g,i}^k\|_2
\le c\Big(\bar{\epsilon}_{\mathrm{cri}}(k)+(k+1)^{-2}+T_k\beta_{k-T_k+1}+\xi_{k+1}\Big).
\label{eq:appD_og_bound}
\end{equation}
Here $\bar{\epsilon}_{\mathrm{cri}}(k)\triangleq \max_{0\le j\le I}\epsilon_{k,j}^{\mathrm{cri}}$.

Then $\mathbb{E}[\Xi_{k+1}]$ can be expanded as: 
\begin{align}
\mathbb{E}[\Xi_{k+1}]
&=\mathbb{E}\Big[\big\langle \boldsymbol{y}_k,\mathbb{E}[\tilde{\boldsymbol{g}}_i^{k+1}\mid \mathcal{F}_k]-\nabla_{\boldsymbol{\theta}}J_i(\boldsymbol{\theta}_k)\big\rangle\Big] \notag\\
&=\mathbb{E}\Big[\big\langle \boldsymbol{y}_k,\boldsymbol{o}_{g,i}^k\big\rangle\Big]
+\mathbb{E}\Big[\big\langle \boldsymbol{y}_k,\nabla_{\boldsymbol{\theta}}\hat{J}_i(\boldsymbol{\theta}_k)-\nabla_{\boldsymbol{\theta}}J_i(\boldsymbol{\theta}_k)\big\rangle\Big].
\label{eq:appD_Xi_expand}
\end{align}
We bound the two terms on the right-hand side separately.

For the first term, $\|\boldsymbol{y}_k\|_2$ is uniformly bounded in expectation because both $\hat{\boldsymbol{g}}_i^k$ and $\nabla_{\boldsymbol{\theta}}J_i(\boldsymbol{\theta}_k)$ are uniformly bounded. Hence, by \eqref{eq:appD_og_bound},
\begin{equation}
\begin{aligned}
\left|\mathbb{E}\Big[\big\langle \boldsymbol{y}_k,\boldsymbol{o}_{g,i}^k\big\rangle\Big]\right|
&\le c\,\mathbb{E}[\|\boldsymbol{o}_{g,i}^k\|_2]\\
&\le c\Big(\bar{\epsilon}_{\mathrm{cri}}(k)+(k+1)^{-2}\\
&\quad+T_k\beta_{k-T_k+1}+\xi_{k+1}\Big).
\end{aligned}
\label{eq:appD_Xi_first}
\end{equation}

Second, by \eqref{eq:appC_grad_diff_reduce}--\eqref{eq:appC_Q_mismatch_bound},
\begin{equation}
\big\|\nabla_{\boldsymbol{\theta}}\hat{J}_i(\boldsymbol{\theta}_k)-\nabla_{\boldsymbol{\theta}}J_i(\boldsymbol{\theta}_k)\big\|_2
\le c\Big(\big|\hat{J}_i^k-J_i(\boldsymbol{\theta}_k)\big|+\epsilon_{m_Q}\Big).
\label{eq:appD_gradgap}
\end{equation}
Again using the boundedness of $\|\boldsymbol{y}_k\|_2$,
\begin{align}
\left|\mathbb{E}\Big[\big\langle \boldsymbol{y}_k,\nabla_{\boldsymbol{\theta}}\hat{J}_i(\boldsymbol{\theta}_k)-\nabla_{\boldsymbol{\theta}}J_i(\boldsymbol{\theta}_k)\big\rangle\Big]\right|
&\le c\,\mathbb{E}\Big[\big|\hat{J}_i^k-J_i(\boldsymbol{\theta}_k)\big|\Big]+c\epsilon_{m_Q}.
\label{eq:appD_Xi_second_pre}
\end{align}
Then we can obtain: 
\begin{align}
\frac{1}{t+1}\sum_{k=0}^{t}\mathbb{E}\Big[\big|\hat{J}_i^k-J_i(\boldsymbol{\theta}_k)\big|\Big]
&\le \sqrt{\frac{1}{t+1}\sum_{k=0}^{t}\mathbb{E}\Big[\big|\hat{J}_i^k-J_i(\boldsymbol{\theta}_k)\big|^2\Big]} \notag\\
&\le \sqrt{\epsilon_J(t)}. 
\label{eq:appD_abs_to_epsJ}
\end{align}
Combining \eqref{eq:appD_Xi_expand}--\eqref{eq:appD_abs_to_epsJ}, we conclude that
\begin{align}
M_4(t)
&\le \mathcal{O}\!\Bigg(\frac{1}{t+1}\sum_{k=0}^{t}\Big(\bar{\epsilon}_{\mathrm{cri}}(k)+(k+1)^{-2}+T_k\beta_{k-T_k+1}+\xi_{k+1}\Big)\Bigg) \notag\\
&\quad +\mathcal{O}\!\Big(\sqrt{\epsilon_J(t)}+\epsilon_{m_Q}\Big).
\label{eq:appD_M4_bound}
\end{align}

For $M_5(t)$, utilizing Cauchy-Schwarz inequality, we have: 
\begin{align}
|M_5(t)|
&\le \frac{1}{t+1}\sum_{k=0}^{t}\frac{1}{\alpha_{k+1}}\mathbb{E}\big[\|\boldsymbol{y}_k\|_2\,\|\boldsymbol{\Delta}_k\|_2\big] \notag\\
&\le \sqrt{\frac{1}{t+1}\sum_{k=0}^{t}\mathbb{E}[\|\boldsymbol{y}_k\|_2^2]}
\sqrt{\frac{1}{t+1}\sum_{k=0}^{t}\frac{\mathbb{E}[\|\boldsymbol{\Delta}_k\|_2^2]}{\alpha_{k+1}^2}}.
\label{eq:appD_M5_pre}
\end{align}
With the bound on $\mathbb{E}[\|\boldsymbol{\Delta}_k\|_2^2]$ derived in $M_2(t)$, 
\begin{equation}
\frac{1}{t+1}\sum_{k=0}^{t}\frac{\mathbb{E}[\|\boldsymbol{\Delta}_k\|_2^2]}{\alpha_{k+1}^2}
\le c\,\frac{1}{t+1}\sum_{k=0}^{t}\frac{\beta_k^2}{\alpha_{k+1}^2}.
\end{equation}
Therefore,
\begin{equation}
\begin{aligned}
M_5(t)
&=\mathcal{O}\!\Big(\sqrt{\epsilon_{g,i}(t)}\,\sqrt{N(t)}\Big),\\
N(t)
&\triangleq \frac{1}{t+1}\sum_{k=0}^{t}\frac{\beta_k^2}{\alpha_{k+1}^2}.
\end{aligned}
\label{eq:appD_M5_bound}
\end{equation}

Substituting \eqref{eq:appD_M1_bound}, \eqref{eq:appD_M2_bound}, \eqref{eq:appD_M3_bound}, \eqref{eq:appD_M4_bound}, and \eqref{eq:appD_M5_bound} into \eqref{eq:appD_gradient_decomp}, we obtain: 
\begin{equation}
\begin{aligned}
\epsilon_{g,i}(t)
&\le \mathcal{O}\!\Big(\sqrt{\epsilon_{g,i}(t)}\,\sqrt{N(t)}\Big)
+\mathcal{O}\!\Big(\frac{1}{(t+1)\alpha_{t+1}}\Big)\\
&\quad +\mathcal{O}\!\Bigg(\frac{1}{t+1}\sum_{k=0}^{t}\Big(
\bar{\epsilon}_{\mathrm{cri}}(k)+(k+1)^{-2}+T_k\beta_{k-T_k+1}+\alpha_{k+1}
\Big)\Bigg)\\
&\quad +\mathcal{O}\!\Bigg(\frac{1}{t+1}\sum_{k=0}^{t}
\frac{\beta_k^2}{\alpha_{k+1}}\Bigg)
+\mathcal{O}\!\Bigg(\frac{1}{t+1}\sum_{k=0}^{t}\xi_{k+1}\Bigg)\\
&+\mathcal{O}\!\Big(\sqrt{\epsilon_J(t)}+\epsilon_{m_Q}\Big).
\end{aligned}
\label{eq:appD_epsgi_pre_solve}
\end{equation}
Let
\begin{equation}
F_i(t)\triangleq \epsilon_{g,i}(t)
\end{equation}
and collect all terms other than $\sqrt{F_i(t)}\sqrt{N(t)}$ into
\begin{equation}
\begin{aligned}
R(t)
&\triangleq \mathcal{O}\!\Big(\frac{1}{(t+1)\alpha_{t+1}}\Big)
+\mathcal{O}\!\Bigg(\frac{1}{t+1}\sum_{k=0}^{t}\Big(
\bar{\epsilon}_{\mathrm{cri}}(k)+(k+1)^{-2}+T_k\beta_{k-T_k+1}
\\&+\alpha_{k+1}\Big)\Bigg)+\mathcal{O}\!\Bigg(\frac{1}{t+1}\sum_{k=0}^{t}
\frac{\beta_k^2}{\alpha_{k+1}}\Bigg)
+\mathcal{O}\!\Bigg(\frac{1}{t+1}\sum_{k=0}^{t}\xi_{k+1}\Bigg)\\
&\quad +\mathcal{O}\!\Big(\sqrt{\epsilon_J(t)}+\epsilon_{m_Q}\Big).
\end{aligned}
\label{eq:appD_R_def}
\end{equation}
Then \eqref{eq:appD_epsgi_pre_solve} takes the form
\begin{equation}
F_i(t)\le c_1\sqrt{F_i(t)}\sqrt{N(t)}+c_2R(t)
\label{eq:appD_quadratic_ineq}
\end{equation}
for some constants $c_1,c_2>0$. Set $u\triangleq \sqrt{F_i(t)}\ge 0$. Then solving this quadratic inequality, we obtain: 
\begin{equation}
\begin{aligned}
u
&\le c_1\sqrt{N(t)}+\sqrt{c_2R(t)},\\
\epsilon_{g,i}(t)
&=\mathcal{O}\!\big(N(t)+R(t)\big).
\end{aligned}
\label{eq:appD_epsgi_intermediate}
\end{equation}
Finally, substituting the definition of $N(t)$ and $R(t)$ into \eqref{eq:appD_epsgi_intermediate}, we obtain:

\begin{equation}
\begin{aligned}
\epsilon_{g,i}(t)
&\le \mathcal{O}\!\Big(\frac{1}{(t+1)\alpha_{t+1}}\Big)\\
&\quad +\mathcal{O}\!\Bigg(\frac{1}{t+1}\sum_{k=0}^{t}\Big(
\bar{\epsilon}_{\mathrm{cri}}(k)+(k+1)^{-1}+T_k\beta_{k-T_k+1}\Big)\Bigg)\\
&\quad +\mathcal{O}\!\Bigg(\frac{1}{t+1}\sum_{k=0}^{t}\Big(
\beta_k^2(\alpha_{k+1}^{-2}+\alpha_{k+1}^{-1})+\alpha_{k+1}\Big)\Bigg)\\
&\quad +\mathcal{O}\!\Bigg(\frac{1}{t+1}\sum_{k=0}^{t}\xi_{k+1}\Bigg)
{}+\mathcal{O}\!\Big(\sqrt{\epsilon_J(t)}+\epsilon_{m_Q}\Big).
\end{aligned}
\label{eq:appD_epsgi_final_fixed_i}
\end{equation}

Then the accumulated gradient-tracking error $\epsilon_g(t)$ can be bounded by the maximum of $\epsilon_{g,i}(t)$ over $i$:
\begin{equation}
\epsilon_g(t)=\frac{1}{t+1}\sum_{k=0}^{t}\max_{0\le j\le I}\mathbb{E}\Big[\big\|\hat{\boldsymbol{g}}_j^k-\nabla_{\boldsymbol{\theta}}J_j(\boldsymbol{\theta}_k)\big\|_2^2\Big]
\le \sum_{j=0}^{I}\epsilon_{g,j}(t).
\end{equation}
This completes the proof of Lemma~\ref{lemma-convergence-rate-surrogate}. 



\bibliographystyle{IEEEtran}
\bibliography{RLreferences}

\end{document}
